 
%%% DOCUMENTCLASS 
%%%-------------------------------------------------------------------------------

\documentclass[
a4paper, % Stock and paper size.
10pt, % Type size.
% article,
% oneside, 
onecolumn, % Only one column of text on a page.
% openright, % Each chapter will start on a recto page.
% openleft, % Each chapter will start on a verso page.
openany, % A chapter may start on either a recto or verso page.
]{memoir}

%%% PACKAGES 
%%%------------------------------------------------------------------------------

\usepackage[utf8x]{inputenc} % If utf8 encoding
% \usepackage[lantin1]{inputenc} % If not utf8 encoding, then this is probably the way to go
\usepackage[T1]{fontenc}    %
\usepackage[english]{babel} % English please
\usepackage[final]{microtype} % Less badboxes
% \usepackage{kpfonts} %Font
\usepackage{amsmath,amssymb,mathtools} % Math
% \usepackage{tikz} % Figures
\usepackage{graphicx} % Include figures
\usepackage{listings}    % Add code snippets
\usepackage[usenames,dvipsnames,svgnames,table]{xcolor}
\usepackage[linktocpage=true]{hyperref} % Extensive support for hypertext in LaTeX
\usepackage[usenames,dvipsnames]{xcolor} %used for font color
\usepackage{dirtree} % Make directory trees easily
\usepackage{framed} % shade in a region of the page
\usepackage{multirow}  % use rows spanning multiple columns in tables
\usepackage{float}
\usepackage[all]{background}
\usepackage[most]{tcolorbox}
\usepackage{hyperref}

%%% PAGE LAYOUT 
%%%------------------------------------------------------------------------------

\setlrmarginsandblock{0.15\paperwidth}{*}{1} % Left and right margin
\setulmarginsandblock{0.2\paperwidth}{*}{1}  % Upper and lower margin
\checkandfixthelayout

%%% SECTIONAL DIVISIONS
%%%------------------------------------------------------------------------------

\maxsecnumdepth{subsubsection} % Subsections (and higher) are numbered
\setsecnumdepth{subsubsection}

\makeatletter %
\makechapterstyle{standard}{
  \setlength{\beforechapskip}{0\baselineskip}
  \setlength{\midchapskip}{1\baselineskip}
  \setlength{\afterchapskip}{2\baselineskip}
  \renewcommand{\chapterheadstart}{\vspace*{\beforechapskip}}
  \renewcommand{\chapnamefont}{\centering\normalfont\huge}
  \renewcommand{\printchaptername}{\chapnamefont \@chapapp}
  \renewcommand{\chapternamenum}{\space}
  \renewcommand{\chapnumfont}{\normalfont\huge}
  \renewcommand{\printchapternum}{\chapnumfont \thechapter}
  \renewcommand{\afterchapternum}{\par\nobreak\vskip \midchapskip}
  \renewcommand{\printchapternonum}{\vspace*{\midchapskip}\vspace*{5mm}}
  \renewcommand{\chaptitlefont}{\centering\bfseries\HUGE}
  \renewcommand{\printchaptertitle}[1]{\chaptitlefont ##1}
  \renewcommand{\afterchaptertitle}{\par\nobreak\vskip \afterchapskip}
}
\makeatother

\chapterstyle{standard}

\setsecheadstyle{\normalfont\huge\bfseries}
\setsubsecheadstyle{\normalfont\large\bfseries}
\setparaheadstyle{\normalfont\normalsize\bfseries}
\setparaindent{0pt}\setafterparaskip{0pt}

%%% FLOATS AND CAPTIONS
%%%------------------------------------------------------------------------------

\makeatletter                  % You do not need to write [htpb] all the time
\renewcommand\fps@figure{htbp} %
\renewcommand\fps@table{htbp}  %
\makeatother                   %

\captiondelim{\space } % A space between caption name and text
\captionnamefont{\small\bfseries} % Font of the caption name
\captiontitlefont{\small\normalfont} % Font of the caption text

\changecaptionwidth          % Change the width of the caption
\captionwidth{1\textwidth} %

%%% ABSTRACT
%%%------------------------------------------------------------------------------

\renewcommand{\abstractnamefont}{\normalfont\small\bfseries} % Font of abstract title
\setlength{\absleftindent}{0.1\textwidth} % Width of abstract
\setlength{\absrightindent}{\absleftindent}

%%% HEADER AND FOOTER 
%%%------------------------------------------------------------------------------

\makepagestyle{standard} % Make standard pagestyle

\makeatletter                 % Define standard pagestyle
\makeevenfoot{standard}{}{}{} %
\makeoddfoot{standard}{}{}{}  %
\makeevenhead{standard}{\bfseries\thepage\normalfont\qquad\small\leftmark}{}{}
\makeoddhead{standard}{}{}{\small\rightmark\qquad\bfseries\thepage}
% \makeheadrule{standard}{\textwidth}{\normalrulethickness}
\makeatother                  %

\makeatletter
\makepsmarks{standard}{
\createmark{chapter}{both}{shownumber}{\@chapapp\ }{ \quad }
\createmark{section}{right}{shownumber}{}{ \quad }
\createplainmark{toc}{both}{\contentsname}
\createplainmark{lof}{both}{\listfigurename}
\createplainmark{lot}{both}{\listtablename}
\createplainmark{bib}{both}{\bibname}
\createplainmark{index}{both}{\indexname}
\createplainmark{glossary}{both}{\glossaryname}
}
\makeatother                               %

\makepagestyle{chap} % Make new chapter pagestyle

\makeatletter
\makeevenfoot{chap}{}{\small\bfseries\thepage}{} % Define new chapter pagestyle
\makeoddfoot{chap}{}{\small\bfseries\thepage}{}  %
\makeevenhead{chap}{}{}{}   %
\makeoddhead{chap}{}{}{}    %
% \makeheadrule{chap}{\textwidth}{\normalrulethickness}
\makeatother

\nouppercaseheads
\pagestyle{standard}               % Choosing pagestyle and chapter pagestyle
\aliaspagestyle{chapter}{chap} %

%%% NEW COMMANDS
%%%------------------------------------------------------------------------------
% custom commands defined in here (must be called before constants)
%%%%%%%%%%%%%%%%%%%%%%%%%%%%%%%%%%%%%%%%%%%%%%%%%%%%%%%%
%%
%%  formatting constants
%%
%%%%%%%%%%%%%%%%%%%%%%%%%%%%%%%%%%%%%%%%%%%%%%%%%%%%%%%%

% formatting the named variables (from user setup)
\newcommand{\definevariable}[1]{\textcolor{blue}{\{#1\}}}
\newcommand{\definevariablecmd}[1]{\textcolor{cyan}{\{#1\}}}

% --------------------------------------------------------
% formatting the named keywords 
\newcommand{\definekeyword}[1]{\textcolor{red}{\{#1\}}}

% --------------------------------------------------------
% Program should be in small caps
\newcommand{\Program}[1]{\textsc{#1}}

% --------------------------------------------------------
% Double underscore
\newlength\dunder
\settowidth{\dunder}{\_}
\newcommand{\twound}{\rule{1.5\dunder}{0.4pt}}

% --------------------------------------------------------
% formatting for the directories trees
\newcommand{\customdirtree}[1]{
\renewcommand*\DTstylecomment{\ttfamily\textcolor{black}}
\renewcommand*\DTstyle{\ttfamily\textcolor{blue}}
\dirtree{%
#1}
}

% --------------------------------------------------------
% Parameters
% \Parameter{Name}{Description}{variable}{default value}{used in}{location}
\ifdevguide
	\newcommand{\ParameterEntry}[7]{
	\textbf{#1}

	\begin{thighlight}
	\textcolor{brown}{#2} 

	\begin{tabular}{>{\color{red}}l c l}
	&&\\
	#3 & = & #4 \\
	\textcolor{blue}{Used in:}  & \multicolumn{2}{p{10cm}}{#5} \\
	\textcolor{blue}{Defined in:} & \multicolumn{2}{p{10cm}}{#6} \\
	\textcolor{blue}{Level:} & \multicolumn{2}{p{10cm}}{#7} \\
	\end{tabular}
	\end{thighlight}
	}
\else
	\newcommand{\ParameterEntry}[7]{
	\textbf{#1}

	\begin{thighlight}
	\textcolor{brown}{#2} 

	\begin{tabular}{>{\color{red}}l c l}
	&&\\
	#3 & = & #4 \\
	\textcolor{blue}{Used in:}  & \multicolumn{2}{p{10cm}}{#5} \\
	\textcolor{blue}{Defined in:} & \multicolumn{2}{p{10cm}}{#6} \\
	\end{tabular}
	\end{thighlight}
	}
\fi


% --------------------------------------------------------
% define dev sections notes etc

\newcommand{\devsection}[2]{
\ifdevguide
\section{#1}

#2
\fi
}


\newcommand{\DevNote}[1]{
\ifdevguide
\begin{note}
#1
\end{note}
\fi
}
% get code highlighting parameters
\definecolor{codegreen}{rgb}{0,0.6,0}
\definecolor{codegray}{rgb}{0.5,0.5,0.5}
\definecolor{codepurple}{rgb}{0.58,0,0.82}
\definecolor{darkgray}{rgb}{0.18,0.31,0.31}
\definecolor{backcolour1}{rgb}{0.95,0.95,0.92}
\definecolor{backcolour2}{rgb}{1,0.977,0.801}
\definecolor{backcolour3}{rgb}{0.79,0.88,1}

\lstdefinestyle{text}{
    commentstyle=\color{codegreen},
    keywordstyle=\color{magenta},
    numberstyle=\tiny\color{codegray},
    stringstyle=\color{codepurple},
    basicstyle=\linespread{1.2}\ttfamily\footnotesize,
    breakatwhitespace=false,         
    breaklines=true,                 
    captionpos=b,                    
    keepspaces=true,                 
    % numbers=left,                    
    % numbersep=5pt,                  
    showspaces=false,                
    showstringspaces=false,
    showtabs=false,                  
    tabsize=2,
    moredelim=**[is][\color{red}]{@}{@},
    moredelim=**[is][\color{OliveGreen}]{<}{>},
    extendedchars=false,
    columns=fullflexible,
    escapeinside={(*}{*)},%
	  literate={\\@}{{\unichar{"0040}}}1
			 {\\<}{{\unichar{"003C}}}1 
			 {\\>}{{\unichar{"003E}}}1
}

\lstdefinestyle{bashstyle}{
	language=bash,  
    commentstyle=\color{codegreen},
    keywordstyle=\color{magenta},
    numberstyle=\tiny\color{codegray},
    stringstyle=\color{codepurple},
    basicstyle=\linespread{1.2}\ttfamily\footnotesize,
    breakatwhitespace=false,         
    breaklines=true,                 
    captionpos=b,                    
    keepspaces=true,                 
    % numbers=left,                    
    % numbersep=5pt,                  
    showspaces=false,                
    showstringspaces=false,
    showtabs=false,                  
    tabsize=2,
    moredelim=**[is][\color{red}]{@}{@},
    extendedchars=false,
    columns=fullflexible,
    escapeinside={(*}{*)},%
    literate={\\@}{{\unichar{"0040}}}1,
}


\lstdefinestyle{cmdstyle}{
    commentstyle=\color{codegreen},
    keywordstyle=\color{magenta},
    numberstyle=\tiny\color{codegray},
    stringstyle=\color{codepurple},
    basicstyle=\linespread{1.2}\ttfamily\footnotesize\color{white},
    breakatwhitespace=false,         
    breaklines=true,                 
    captionpos=b,                    
    keepspaces=true,                 
    % numbers=left,                    
    % numbersep=5pt,                  
    showspaces=false,                
    showstringspaces=false,
    showtabs=false,                  
    tabsize=2,
    extendedchars=false,
    columns=fullflexible,
    escapeinside={(*}{*)},%
}

\lstdefinestyle{pythonstyle}{
    language=Python,
    commentstyle=\color{codegreen},
    keywordstyle=\color{magenta},
    numberstyle=\tiny\color{codegray},
    stringstyle=\color{codepurple},
    basicstyle=\linespread{1.2}\ttfamily\footnotesize,
    breakatwhitespace=false,         
    breaklines=true,                 
    captionpos=b,                    
    keepspaces=true,                 
    % numbers=left,                    
    % numbersep=5pt,                  
    showspaces=false,                
    showstringspaces=false,
    showtabs=false,                  
    tabsize=2,
    moredelim=**[is][\color{red}]{@}{@},
    extendedchars=false,
    columns=fullflexible,
    escapeinside={(*}{*)},%
    literate={\\@}{{\unichar{"0040}}}1
}



\tcbuselibrary{listings}


\newtcblisting{textbox}{
      colback=backcolour1,
      colframe=backcolour1!75!black,
      listing only,
      listing options={style=text},
}

\newtcblisting{bashbox}{
      colback=darkgray,
      colframe=darkgray!75!black,,
      listing only,
      listing options={style=bashstyle},
}


\newtcblisting{cmdbox}{
      colback=darkgray,
      colframe=darkgray!75!black,,
      listing only,
      listing options={style=cmdstyle},
}

\newtcblisting{pythonbox}{
      colback=backcolour3,
      colframe=backcolour3!75!black,
      listing only,
      listing options={style=pythonstyle},
}
% custom constants defined in here
% ----------------------------------------------
% Program constants that can be changed through-out the document

% ----------------------------------------------
% Manual version, authors and dates
% ----------------------------------------------
% the current versions of the manuals
\newcommand{\MyVersion}{0.0.1}
\newcommand{\MyVersionUser}{\MyVersion}
\newcommand{\MyVersionDev}{\MyVersion}

% the last date modified for the manuals
\newcommand{\MyDate}{2017-11-27}
\newcommand{\MyDateUser}{\MyDate}
\newcommand{\MyDateDev}{\MyDate}

% the author list of the manuals
\newcommand{\MyAuthors}{N. Cook, F. Bouchy, E. Artigau, I. Boisse, M. Hobson, C. Moutou}

% the id of the manuals
\newcommand{\MyID}{SPIROU-4800-LAM-UM-00961}
\newcommand{\MyIDuser}{\MyID}
\newcommand{\MyIDdev}{\MyID}
% ----------------------------------------------
% DRS variables
% ----------------------------------------------
% The version of the DRS code
\newcommand{\MyCodeVersion}{0.0.1}

% The instrument
\newcommand{\instrument}{SPIRou}

% The latest download links
\newcommand{\DRSLatestURL}{https://github.com/njcuk9999/spirou\_py3}

\newcommand{\DRSGitURL}{https://github.com/njcuk9999/spirou\_py3.git}

\newcommand{\DRSsshURL}{git@github.com:njcuk9999/spirou\_py3.git}

% ----------------------------------------------
% Commonly used variables
% ----------------------------------------------

% variable the user should identify with their installation directory
\newcommand{\InstallDIR}{\definevariable{INSTALL\_DIR}}

% name for mac operating system
\newcommand{\mac}{macOS\,}
% Not mine
\newcommand{\p}{\partial} %Partial
% Or what ever you want


%%% MARGIN
%%%------------------------------------------------------------------------------
\SetBgContents{\MyIDdev}% Set contents
\SetBgPosition{-0.75cm,-0.3\textheight}% Select location
\SetBgOpacity{1.0}% Select opacity
\SetBgAngle{90.0}% Select rotation of logo
\SetBgScale{1.25}% Select scale factor of logo

%%% TABLE OF CONTENTS
%%%------------------------------------------------------------------------------

\maxtocdepth{subsection} % Only parts, chapters and sections in the table of contents
\settocdepth{subsection}

% \maxtocdepth{section} % Only parts, chapters and sections in the table of contents
% \settocdepth{section}

\AtEndDocument{\addtocontents{toc}{\par}} % Add a \par to the end of the TOC

\renewcommand*{\cftappendixname}{Appendix\space}

%%% INTERNAL HYPERLINKS
%%%-------------------------------------------------------------------------------

\usepackage{hyperref}   % Internal hyperlinks
\hypersetup{
    pdfauthor={Neil James Cook},
    pdfcreator={Neil James Cook},
    pdftitle={SPIRou Developer Guide},
    pdfsubject={SPIRou DRS},
    pdfkeywords={SPIRou, DRS, Pipeline},
    colorlinks=true,         % false: boxed links; true: colored links
    linkcolor=blue,          % color of internal links (change box color with linkbordercolor)
    citecolor=Maroon,        % color of links to bibliography
    filecolor=blue,          % color of file links
    urlcolor=blue,           % color of external links
    plainpages=false,       
}
\usepackage{memhfixc}   %



%%% THE DOCUMENT
%%% Where all the important stuff is included!
%%%-------------------------------------------------------------------------------

\author{\MyAuthors}
\title{{\Huge SPIRou Data Reduction Software} \vspace{1cm} \\ \HUGE{Developer Guide} \\ {\small \MyVersionDev} \\ For DRS \instrument \MyCodeVersion}
\date{\MyDateDev}

% \usepackage{lipsum} % Just to put in some text

\begin{document}

\frontmatter

\maketitle

% Add logo to title page
\vspace{1cm}
\begin{center}
\includegraphics[width=.5\textwidth]{Figures/Logo_SPIRou-22.pdf}
\end{center}
\vspace{1cm}

\begin{abstract}
\noindent This is the guide to coding the DRS (including installation, running, rules and stardisation approaches). This document is not intended for the general used of the DRS, instead it is intended for those who wish to develop the software further and understand the changes between this version and previous versions.
\end{abstract}
\clearpage

\tableofcontents*
\clearpage

% Introduction
%%%%%%%%%%%%%%%%%%%%%%%%%%%%%%%%%%%%%%%%%%%%%%%%%%%%%%%%
%%
\chapter{Introduction}
\label{chapter:introduction}
%%
%%%%%%%%%%%%%%%%%%%%%%%%%%%%%%%%%%%%%%%%%%%%%%%%%%%%%%%%


%%%-------------------------------------------------------------------------------
\mainmatter
% Define the chapters here
%%%-------------------------------------------------------------------------------

% Chapter 1: Installation guide
%%%%%%%%%%%%%%%%%%%%%%%%%%%%%%%%%%%%%%%%%%%%%%%%%%%%%%%%
%%
\chapter{Installation}
\label{chapter:installation}
%%
%%%%%%%%%%%%%%%%%%%%%%%%%%%%%%%%%%%%%%%%%%%%%%%%%%%%%%%%

%%%%%%%%%%%%%%%%%%%%%%%%%%%%%%%%%%%%%%%%%%%%%%%%%%%%%%%%
%%
\section{Introduction}
\label{ch:install:installintro}
%%
%%%%%%%%%%%%%%%%%%%%%%%%%%%%%%%%%%%%%%%%%%%%%%%%%%%%%%%%

Once finialised the installation should just be a download, run setup.py and configure the DRS directories, however, during development the following stages are required.

\NOTE{Currently the download repositry on github is private and requires a github account, and the user to be added to the list of collaborators. To be added to the collaborators please email \url{neil.james.cook@gmail.com} with your github username.}







%%%%%%%%%%%%%%%%%%%%%%%%%%%%%%%%%%%%%%%%%%%%%%%%%%%%%%%%
%%
\section{Download}
\label{ch:install:installDownload}
%%
%%%%%%%%%%%%%%%%%%%%%%%%%%%%%%%%%%%%%%%%%%%%%%%%%%%%%%%%

Get the latest version of the DRS (for \instrument version \MyCodeVersion). Use any of the following ways:

\begin{itemize}
\item manually download from here: \url{\DRSLatestURL}

\item use Git:
\begin{cmdbox}
git checkout (*\DRSGitURL*)
\end{cmdbox}

\item use SVN:
\begin{cmdbox}
svn checkout  (*\DRSGitURL*)
\end{cmdbox}

\item use ssh:
\begin{cmdbox}
scp -r (*\DRSsshURL*)
\end{cmdbox}

\end{itemize}


%%%%%%%%%%%%%%%%%%%%%%%%%%%%%%%%%%%%%%%%%%%%%%%%%%%%%%%%
%%
\clearpage
\newpage
\section{Prerequisites}
\label{ch:install:prerequisites}
%%
%%%%%%%%%%%%%%%%%%%%%%%%%%%%%%%%%%%%%%%%%%%%%%%%%%%%%%%%

It is recommended to install the latest version of Anaconda python distribution, available for Windows, \mac and Linux (here: \url{https://www.anaconda.com/download/}). However one can run the DRS on a native python installation. \\

\noindent We recommend python 3 over python 2 for long term continued support (however the latest version of the DRS supports the newest versions of python 2.7). \\

\NOTE{Before installing the DRS you must have one of the following:} \\

\noindent EITHER (RECOMMENDED)
\begin{itemize}
\item A valid version of the Anaconda python distribution (for python2 or python 3)
	\noindent Currently tested version of python are:
	\begin{itemize}
	\item Python 2.7.13 and Anaconda 4.4.0
	\item Python 3.6.3 and Anaconda 5.0.1 --- RECOMMENDED
	\end{itemize}
\end{itemize}
\noindent OR
\begin{itemize}
\item An up-to-date version of python and the following python modules (for either python 2 or python 3)
	\begin{itemize}
	\item Python 3.6
		\begin{itemize}
		\item \Program{astropy} (tested with version 2.0.2)
		\item \Program{matplotlib} (tested with version 2.1.0)
		\item \Program{numpy} (tested with version 1.13.3)
		\item and the following built-in modules (comes with python): \Program{datetime}, \Program{filecmp}, \Program{glob}, \Program{os}, \Program{pkg\_resources}, \Program{shutil}, \Program{sys}, \Program{time}, \Program{warnings}
		\end{itemize}
	\item Python 2.7
		\begin{itemize}
		\item astropy (tested with version 1.3.2)
		\item matplotlib (tested with version 2.0.2)
		\item numpy (tested with version 1.12.1)
		\item and the following built-in modules (comes with python): \Program{\twound{future}\twound}, \Program{collections}, \Program{datetime}, \Program{filecmp}, \Program{glob}, \Program{os}, \Program{pkg\_resources}, \Program{shutil}, \Program{sys}, \Program{time}, \Program{warnings}
		\end{itemize}
	\end{itemize}
\end{itemize}



%%%%%%%%%%%%%%%%%%%%%%%%%%%%%%%%%%%%%%%%%%%%%%%%%%%%%%%%
%%
\clearpage
\newpage
\section{Installation Linux and \mac}
\label{ch:install:installunix}
%%
%%%%%%%%%%%%%%%%%%%%%%%%%%%%%%%%%%%%%%%%%%%%%%%%%%%%%%%%

Currently the DRS has to be installed manually. This involves the following steps:
\begin{enumerate}
\item Extraction (Section \ref{ch:install:installunix:extraction})
\item Modify environmental settings (Section \ref{ch:install:installunix:environ_settings})
\item Make recipes executable (Section \ref{ch:install:installunix:executable})
\end{enumerate}

% -------------------------------------------------------
\subsection{Extraction}
\label{ch:install:installunix:extraction}
% -------------------------------------------------------

The first step is to extract the DRS into a folder (the \InstallDIR).

\noindent Do this by using the following commands:
\begin{cmdbox}
cd INSTALL_DIR 
unzip DRS.zip
\end{cmdbox}

% -------------------------------------------------------
\subsection{Modify environmental settings}
\label{ch:install:installunix:environ_settings}
% -------------------------------------------------------

The next step is to modify your PATH and PYTHONPATH environmental variables (to include the \InstallDIR. This depends which shell you are using (type `\lstinline[style=bashstyle]{echo $0}' to find out which).

\begin{itemize}
	\item In bash open the `.bashrc' text file in your home ($\sim$) directory (or create it if it doesn't exist)

	\begin{textbox}
	@export@ PATH=(*\InstallDIR*)/bin/:$PATH

	@export@ PYTHONPATH=(*\InstallDIR*):(*\InstallDIR*)/bin/:$PYTHONPATH
	\end{textbox}

	\item In csh /tcsh open the `.cshrc' or `.tcshrc' text file in your home ($\sim$) directory (or create it if it doesn't exist) 

	\begin{textbox}
	@setenv@ PATH (*\InstallDIR*)/bin/:@${PATH}@

	@setenv@ PYTHONPATH (*\InstallDIR*):(*\InstallDIR*)/bin/:${PYTHONPATH}
	\end{textbox}

\end{itemize}

% -------------------------------------------------------
\subsection{Make recipes executable}
\label{ch:install:installunix:executable}
% -------------------------------------------------------

\noindent To run the recipes from the command line (without starting python) one must make them executable. Do this by using the following command:
\begin{bashbox}
chmod +x (*\InstallDIR*)/bin/*.py
\end{bashbox}



%%%%%%%%%%%%%%%%%%%%%%%%%%%%%%%%%%%%%%%%%%%%%%%%%%%%%%%%
%%
\clearpage
\newpage
\section{Installation Windows}
\label{ch:install:install_win}
%%
%%%%%%%%%%%%%%%%%%%%%%%%%%%%%%%%%%%%%%%%%%%%%%%%%%%%%%%%

This is very similar currently to the Linux/\mac installation (in the future a `.exe' file will be given).

\begin{enumerate}
\item Extract to \InstallDIR with your favourite unzipping softwear.
\item Add \InstallDIR to your PYTHONPATH (Section \ref{ch:install:install_win:environ_settings})
\end{enumerate}

% -------------------------------------------------------
\subsection{How to modify environmental settings in windows}
\label{ch:install:install_win:environ_settings}
% -------------------------------------------------------

This process is a little more convoluted than on Linux or \mac system.

\begin{enumerate}
\item Go to `My computer > Properties > Advanced System Settings > Enviromental Variables'.

\item if under system variable `PythonPath' exists click edit and add `\InstallDIR;'\, to the end.

\noindent i.e.

\begin{textbox}
C:\Python27;(*\InstallDIR;*)
\end{textbox}

\item if under system variables `PythonPath' does not exist create a new variable called `PythonPath' and add:

\begin{textbox}
%PYTHONPATH%;(*\InstallDIR*);(*\InstallDIR*)\bin\;
\end{textbox}

\end{enumerate}

\noindent For problems/troubleshooting see here: \url{https://stackoverflow.com/questions/3701646/how-to-add-to-the-pythonpath-in-windows-7}.



%%%%%%%%%%%%%%%%%%%%%%%%%%%%%%%%%%%%%%%%%%%%%%%%%%%%%%%%
%%
\clearpage
\newpage
\section{Setting up the DRS}
\label{ch:install:setup}
%%
%%%%%%%%%%%%%%%%%%%%%%%%%%%%%%%%%%%%%%%%%%%%%%%%%%%%%%%%

Before running the DRS one must set the data paths. \\

% This can be done two ways:

% \noindent EITHER
% \begin{itemize}
% 	\item Run the `DRS\_setup.py' script (Section \ref{ch:install:setup:drs_setup_script}).
% \end{itemize}
% \noindent OR
% \begin{itemize}
% 	\item Manually edit the config.txt file (Section \ref{ch:install:setup:manual_edit_textfile}).
% \end{itemize}


% % -------------------------------------------------------
% \subsection{Running the `DRS\_setup.py' script}
% \label{ch:install:setup:drs_setup_script}
% % -------------------------------------------------------

% DRS\_setup.py script not yet written please use manual editting of config.txt file.

% % -------------------------------------------------------
% \subsection{Manually editting the config.txt file}
% \label{ch:install:setup:manual_edit_textfile}
% % -------------------------------------------------------

\noindent The `config.txt' file is located in the \InstallDIR in the config folder. \\

i.e. at \InstallDIR{/config/config.txt} \\

\noindent The following keywords must be changed (and must be a valid path):

\noindent\fboxsep=0pt%

\highlight{0.6}{1.1}{
\begin{table}[H]
\begin{tabular}{l c l c p{6.25cm}}
\definevariable{TDATA}            & = & /drs/data/        & / & Define the DATA directory\\
\definevariable{DRS\_ROOT}         & = & /drs/INTROOT/     & / & Define the installation directory (\InstallDIR) \\
\definevariable{DRS\_DATA\_RAW}     & = & /drs/data/raw     & / & Define the folder with the raw data files in \\
\definevariable{DRS\_DATA\_REDUC}   & = & /drs/data/reduced & / & Define the directory that the reduced data should be saved to/read from \\
\definevariable{DRS\_CALIB\_DB}     & = & /drs/data/calibDB & / & Define the directory that the calibration files should be saved to/read from \\
\definevariable{DRS\_DATA\_MSG}     & = & /drs/data/msg     & / & Define the directory that the log messages are stored in \\
\definevariable{DRS\_DATA\_WORKING} & = & /drs/data/tmp/    & / & Define the working directory \\
\end{tabular}
\end{table}

\setlength{\leftskip}{0.5 cm}

\noindent The directories here are for linux and \mac systems another example would be `/home/user/INTROOT' for the \InstallDIR directory. On Windows machines this would be equivalent to `C:\textbackslash{Users}\textbackslash{<username>}\textbackslash{INTROOT}' in Windows Vista, 7, 8 and 10 or `C:\textbackslash{Documents and Settings}\textbackslash{<username>}\textbackslash{INTROOT}' on early versions of Windows. \\
} \\

\setlength{\leftskip}{0pt}

\noindent The following keywords can be changed: \\

\highlight{0.4}{1.1}{
\begin{table}[H]
\begin{tabular}{>{\color{red}}l c r c p{5cm}}
\definevariable{DRS\_PLOT}    & = & 1     & / & Whether to show plots \\
\definevariable{PRINT\_LEVEL} & = & "all" & / & Level at which to print \\
\definevariable{LOG\_LEVEL}   & = & "all" & / & Level at which to log in log file \\
\end{tabular}
\end{table}

\noindent For the `\definevariable{PRINT\_LEVEL} and \definevariable{LOG\_LEVEL} keywords the values are set as follows:
\begin{itemize}
	\item "all" -- prints all events
	\item "info" -- prints info, warning and error events
	\item "warning" -- prints warning and error events
	\item "error" -- print only error events
\end{itemize}
}

%%%%%%%%%%%%%%%%%%%%%%%%%%%%%%%%%%%%%%%%%%%%%%%%%%%%%%%%
%%
\clearpage
\newpage
\section{Validating Installion on Linux and macOS}
\label{ch:install:validating_installunix}
%%
%%%%%%%%%%%%%%%%%%%%%%%%%%%%%%%%%%%%%%%%%%%%%%%%%%%%%%%%

Note one must install the DRS (Section \ref{ch:install:installunix}) AND set up the DRS (Section \ref{ch:install:setup}). \\

There are four ways to run the DRS in Linux and \mac (thus four ways to verify installation was correct).

\begin{itemize}

\item To validate running from command line type:
\begin{cmdbox}
cal_validate_drs.py
\end{cmdbox}

\item To validate running from python/ipython from the command line type:
\begin{cmdbox}
python cal_validate_drs.py

ipython cal_validate_drs.py
\end{cmdbox}

\item To validate running from ipython, open ipython and type:
\begin{pythonbox}
@run@ cal_validate_drs.main()
\end{pythonbox}

\item To validate running from import from python/ipython, open python/ipython and type:
\begin{pythonbox}
import cal_validate_drs

cal_validate_drs.main()
\end{pythonbox}

\end{itemize}

\noindent If validation is successful the following should appear:

\begin{cmdbox}
22:11:44.1 -   || *****************************************
22:11:44.1 -   || * SPIROU @(#) Geneva Observatory (0.0.1)
22:11:44.1 -   || *****************************************
22:11:44.1 -   ||(dir_data_raw)      DRS_DATA_RAW=/scratch/Projects/spirou_py3/data/raw
22:11:44.1 -   ||(dir_data_reduc)    DRS_DATA_REDUC=/scratch/Projects/spirou_py3/data/reduced
22:11:44.1 -   ||(dir_calib_db)      DRS_CALIB_DB=/scratch/Projects/spirou_py3/data/calibDB
22:11:44.1 -   ||(dir_data_msg)      DRS_DATA_MSG=/scratch/Projects/spirou_py3/data/msg
22:11:44.1 -   ||(print_level)       PRINT_LEVEL=all         %(error/warning/info/all)
22:11:44.1 -   ||(log_level)         LOG_LEVEL=all         %(error/warning/info/all)
22:11:44.1 -   ||(plot_graph)        DRS_PLOT=1            %(def/undef/trigger)
22:11:44.1 -   ||(used_date)         DRS_USED_DATE=undefined
22:11:44.1 -   ||(working_dir)       DRS_DATA_WORKING=/scratch/Projects/spirou_py3/data/tmp/
22:11:44.1 -   ||                    DRS_INTERACTIVE is not set, running on-line mode
22:11:44.1 -   ||
22:11:44.1 -   ||Validation successful. DRS installed corrected.
\end{cmdbox}

%%%%%%%%%%%%%%%%%%%%%%%%%%%%%%%%%%%%%%%%%%%%%%%%%%%%%%%%
%%
\clearpage
\newpage
\section{Validating Installion on Windows}
\label{ch:install:validating_installwin}
%%
%%%%%%%%%%%%%%%%%%%%%%%%%%%%%%%%%%%%%%%%%%%%%%%%%%%%%%%%

Note one must install the DRS (Section \ref{ch:install:install_win}) AND set up the DRS (Section \ref{ch:install:setup}). \\

In windows there are currently 3 ways to run the RS (running in python/ipython).

\begin{itemize}

\item To validate running from python/ipython from the command line type:
\begin{cmdbox}
python cal_validate_drs.py

ipython cal_validate_drs.py
\end{cmdbox}

\item To validate running from ipython, open ipython and type:
\begin{pythonbox}
@run@ cal_validate_drs.main()
\end{pythonbox}

\item To validate running from import from python/ipython, open python/ipython and type:
\begin{pythonbox}
import cal_validate_drs

cal_validate_drs.main()
\end{pythonbox}

\end{itemize}

\noindent If validation is successful the following should appear:

\begin{cmdbox}
22:11:44.1 -   || *****************************************
22:11:44.1 -   || * SPIROU @(#) Geneva Observatory (0.0.1)
22:11:44.1 -   || *****************************************
22:11:44.1 -   ||(dir_data_raw)      DRS_DATA_RAW=/scratch/Projects/spirou_py3/data/raw
22:11:44.1 -   ||(dir_data_reduc)    DRS_DATA_REDUC=/scratch/Projects/spirou_py3/data/reduced
22:11:44.1 -   ||(dir_calib_db)      DRS_CALIB_DB=/scratch/Projects/spirou_py3/data/calibDB
22:11:44.1 -   ||(dir_data_msg)      DRS_DATA_MSG=/scratch/Projects/spirou_py3/data/msg
22:11:44.1 -   ||(print_level)       PRINT_LEVEL=all         %(error/warning/info/all)
22:11:44.1 -   ||(log_level)         LOG_LEVEL=all         %(error/warning/info/all)
22:11:44.1 -   ||(plot_graph)        DRS_PLOT=1            %(def/undef/trigger)
22:11:44.1 -   ||(used_date)         DRS_USED_DATE=undefined
22:11:44.1 -   ||(working_dir)       DRS_DATA_WORKING=/scratch/Projects/spirou_py3/data/tmp/
22:11:44.1 -   ||                    DRS_INTERACTIVE is not set, running on-line mode
22:11:44.1 -   ||
22:11:44.1 -   ||Validation successful. DRS installed corrected.
\end{cmdbox}

% Chapter 2: Data Architecture
%%%%%%%%%%%%%%%%%%%%%%%%%%%%%%%%%%%%%%%%%%%%%%%%%%%%%%%%
%%
\chapter{Data Architecture}
\label{chapter:data_architecture}
%%
%%%%%%%%%%%%%%%%%%%%%%%%%%%%%%%%%%%%%%%%%%%%%%%%%%%%%%%%

% Chapter 3: Using the DRS
%%%%%%%%%%%%%%%%%%%%%%%%%%%%%%%%%%%%%%%%%%%%%%%%%%%%%%%%
%%
\chapter{Using the DRS}
\label{chapter:using_the_drs}
%%
%%%%%%%%%%%%%%%%%%%%%%%%%%%%%%%%%%%%%%%%%%%%%%%%%%%%%%%%

% Chapter 4: Changelog (This version and AT-4)
%%%%%%%%%%%%%%%%%%%%%%%%%%%%%%%%%%%%%%%%%%%%%%%%%%%%%%%%
%%
\chapter{Summary of changes (from version AT-4)}
\label{chapter:the_changelog}
%%
%%%%%%%%%%%%%%%%%%%%%%%%%%%%%%%%%%%%%%%%%%%%%%%%%%%%%%%%

% Chapter 5: Coding rules and standardisation practises
%%%%%%%%%%%%%%%%%%%%%%%%%%%%%%%%%%%%%%%%%%%%%%%%%%%%%%%%
%%
\chapter{Coding rules and standardisation practises}
\label{chapter:rules_and_standardisation}
%%
%%%%%%%%%%%%%%%%%%%%%%%%%%%%%%%%%%%%%%%%%%%%%%%%%%%%%%%%

% Chapter 6: Required input header keywords
%%%%%%%%%%%%%%%%%%%%%%%%%%%%%%%%%%%%%%%%%%%%%%%%%%%%%%%%
%%
\chapter{Required input header keywords}
\label{ch:input_keywords}
%%
%%%%%%%%%%%%%%%%%%%%%%%%%%%%%%%%%%%%%%%%%%%%%%%%%%%%%%%%

%%%%%%%%%%%%%%%%%%%%%%%%%%%%%%%%%%%%%%%%%%%%%%%%%%%%%%%%
%%
\section{Required keywords}
%%
%%%%%%%%%%%%%%%%%%%%%%%%%%%%%%%%%%%%%%%%%%%%%%%%%%%%%%%%

The following keywords are required by the current recipes to run.

\begin{itemize}

% TPL_NAME
\item \KeywordEntry{Data fits file type}
{The data fits file type (template Name)}
{kw\_DPRTYPE}
{TPL\_NAME}{"DATA"}{template Name}
{\AllRecipes}{\spirouKeywords}{\spirouKeywords}

% ACQTIME1
\item \label{text:acqtime1} \KeywordEntry{Acquisition time (human readable)}
{The acquisition time in format YYYY-mm-dd-HH-MM-SS.ss}
{kw\_ACQTIME\_KEY}
{ACQTIME1}{YYYY-mm-dd-HH-MM-SS.ss}{Date at start of observation}
{\AllRecipes}{\spirouKeywords}{\spirouKeywords}

% ACQTIME
\item \KeywordEntry{Acquisition time (unix time format)}
{The acquisition time in in unix time format (time since 1970-01-01-00-00-00)}
{kw\_ACQTIME\_KEY\_UNIX}
{ACQTIME}{000000000.00}{Date in unix time at start of observation}
{\AllRecipes}{\spirouKeywords}{\spirouKeywords}

% RDNOISE
\item \KeywordEntry{Read noise}
{The read noise (used for sigdet) [e-]}
{kw\_RDNOISE}
{RDNOISE}{0.0}{read noise (electrons)}
{\AllRecipes}{\spirouKeywords}{\spirouKeywords}

% GAIN
\item \KeywordEntry{Gain}
{The gain [e-/ADU]}
{kw\_GAIN}
{GAIN}{0.0}{gain (electrons/ADU)}
{\AllRecipes}{\spirouKeywords}{\spirouKeywords}

% EXPTIME
\item \KeywordEntry{Exposure time}
{The integration time in seconds}
{kw\_EXPTIME}
{EXPTIME}{0.0}{Integration time (seconds)}
{\AllRecipes}{\spirouKeywords}{\spirouKeywords}

\end{itemize}

\ifdevguide

%%%%%%%%%%%%%%%%%%%%%%%%%%%%%%%%%%%%%%%%%%%%%%%%%%%%%%%%
%%
\section{Descriptions}
%%
%%%%%%%%%%%%%%%%%%%%%%%%%%%%%%%%%%%%%%%%%%%%%%%%%%%%%%%%

% ------------------------------------------------------------------------

The following FITS descriptors of the 2D raw frames are required for the DRS.
Last updated version 21 Nov 2014. 

%%%%%%%%%%%%%%%%%%%%%%%%%%%%%%%%%%%%%%%%%%%%%%%%%%%%%%%%
%%
\subsection{Standard FITS Keywords}
%%
%%%%%%%%%%%%%%%%%%%%%%%%%%%%%%%%%%%%%%%%%%%%%%%%%%%%%%%%

\begin{thighlight}
\begin{table}[H]
\begin{tabular}{>{\color{red}}l c r c l}
BITPIX  & = &                   16 & / & 16bit \\
NAXIS   & = &                    2 & / & Number of axes \\
NAXIS1  & = &                 4096 & / & Number of pixel columns \\
NAXIS2  & = &                 4096 & / & Number of pixel rows \\
BZERO   & = &              32768.0 & / & Zero factor \\
BSCALE  & = &                  1.0 & / & Scale factor \\
DATE    & = & `2013-11-26T09:06:14' & / & UTC Date of file creation \\
INSTRUME& = & `SPIROU'           & / & Instrument Name \\
\end{tabular}
\end{table}
\end{thighlight}

%%%%%%%%%%%%%%%%%%%%%%%%%%%%%%%%%%%%%%%%%%%%%%%%%%%%%%%%
%%
\subsection{FITS keywords related to the detector}
%%
%%%%%%%%%%%%%%%%%%%%%%%%%%%%%%%%%%%%%%%%%%%%%%%%%%%%%%%%

\begin{thighlight}
\begin{table}[H]
\begin{tabular}{>{\color{red}}l c r c l}
EXPTIME & = &                800.0 & / &  Integration time (seconds) \\
DARKTIME& = &            800.0 & / & Dark current time (seconds) \\
GAIN    & = &                 1.30 & / & gain (electrons/ADU) \\
RDNOISE & = &                 4.20 & / & read noise (electrons) \\
NSUBEXP & = &                    4 & / & Total number of sub-exposures of 5.2s \\
OBSTYPE & = &   `NORMAL'     & / & Exposure type (DARK/NORMAL) \\
MIDEXPTM& = &        400  & / &  mid-exposure time (seconds)  \\
EMCNTS  & = & 	444578   & / & exposure meter counts at end \\
\end{tabular}
\end{table}
\end{thighlight}

%%%%%%%%%%%%%%%%%%%%%%%%%%%%%%%%%%%%%%%%%%%%%%%%%%%%%%%%
%%
\subsection{FITS keywords related to the target}
%%
%%%%%%%%%%%%%%%%%%%%%%%%%%%%%%%%%%%%%%%%%%%%%%%%%%%%%%%%

\begin{thighlight}
\begin{table}[H]
\begin{tabular}{>{\color{red}}l c r c l}
OBJNAME & = &  Gl9999   & / &  Target name \\
OBJRA   & = &  `5:35:09.87'         & / & Target right ascension \\
OBJDEC  & = &  `-5:27:53.3'        & / & Target declination \\
OBJRAPM & = &                  0.560 & / & Target right ascension proper motion in as/yr \\
OBJDECPM& = &                  -0.33 & / & Target declination proper motion in as/yr \\
OBJEQUIN& = &  2000.0       & / & Target equinox \\

OBJRV   & = &        -30.0      & / & Target Radial velocity (km/s)  (999 if unknown) \\
OBJTYPE & = &     `M5' & / & Target spectral type \\
OBJJMAG & = &        8.2 & / & Target J magnitude \\
OBJHMAG & = &        9.2 & / & Target H magnitude \\
OBJKMAG & = &        10.0 & / & Target K magnitude  \\
\end{tabular}
\end{table}
\end{thighlight}

%%%%%%%%%%%%%%%%%%%%%%%%%%%%%%%%%%%%%%%%%%%%%%%%%%%%%%%%
%%
\subsection{FITS keywords related to the telescope}
%%
%%%%%%%%%%%%%%%%%%%%%%%%%%%%%%%%%%%%%%%%%%%%%%%%%%%%%%%%

\begin{thighlight}
\begin{table}[H]
\begin{tabular}{>{\color{red}}l c r c l}
ACQTIME & = & 2013-11-26T09 :06 :14.858 & / & Date at start of observation \\
ACQTIME1 & = & 1385456774               & / & Date in unix time at start of observation \\
DATE\_OBS& = & `2013-11-26T09 :06 :14.858'       & / & Date at start of observation (UTC) \\
EQUINOX & = &               2000.0 & / &Equinox of coordinates \\
EPOCH   & = &            2000.0 & / & Epoch of coordinates \\
MJDATE  & = &        56622.3700212 & / & Modified Julian Date at start of observation \\
MJEND   & = &        56622.3797593 & / & Modified Julian Date at end of observation \\
AIRMASS & = &        1.4 & / & Airmass at start of observation \\
RA      & = & `5:35:09.87'         & / & Telescope right ascension \\
DEC     & = & `-5:27:53.3'         & / & Telescope declination \\

SEEING & = &    1.0 & / & Seeing at start of observation \\
\end{tabular}
\end{table}
\end{thighlight}

%%%%%%%%%%%%%%%%%%%%%%%%%%%%%%%%%%%%%%%%%%%%%%%%%%%%%%%%
%%
\subsection{FITS keywords related to the instrument}
%%
%%%%%%%%%%%%%%%%%%%%%%%%%%%%%%%%%%%%%%%%%%%%%%%%%%%%%%%%

\begin{thighlight}
\begin{table}[H]
\begin{tabular}{>{\color{red}}l c r c l}
TPL\_NAME & = & `SPIROU\_POL\_WAVE'   & / & template Name  \\
TPL\_NEXP & = &    1                  & / & \# of exposure within template  \\
TPL\_EXPN & = &    1                  & / & exposure \# within template  \\
INS\_CAL  & = & `WAVE'                & / & Simultaneous calibration (WAVE/FP/NONE)  \\
INS\_LAMP & = & `UrAr'                & / &   Calibration lamp  \\
INS\_RHB1 & = &  90                   & / & SPIROU rhomb 1 position (deg) \\
INS\_RHB2 & = &  180                  & / & SPIROU rhomb 2 position deg) \\
\end{tabular}
\end{table}
\end{thighlight}

\fi

% Chapter 7: Descriptions of constants and variables
%%%%%%%%%%%%%%%%%%%%%%%%%%%%%%%%%%%%%%%%%%%%%%%%%%%%%%%%
%%
\chapter{Description of Constants and Variables}
\label{chapter:constants_and_variables}
%%
%%%%%%%%%%%%%%%%%%%%%%%%%%%%%%%%%%%%%%%%%%%%%%%%%%%%%%%%

% Chapter 8: The Recipes
%%%%%%%%%%%%%%%%%%%%%%%%%%%%%%%%%%%%%%%%%%%%%%%%%%%%%%%%
%%
\chapter{The Recipes}
\label{ch:the_recipes}
%%
%%%%%%%%%%%%%%%%%%%%%%%%%%%%%%%%%%%%%%%%%%%%%%%%%%%%%%%%

% Chapter 9: The Recipes
%%%%%%%%%%%%%%%%%%%%%%%%%%%%%%%%%%%%%%%%%%%%%%%%%%%%%%%%
%%
\chapter{The DRS Module}
\label{ch:the_module}
%%
%%%%%%%%%%%%%%%%%%%%%%%%%%%%%%%%%%%%%%%%%%%%%%%%%%%%%%%%

\section{Introduction}
\label{ch:the_module:intro}

In the below sections the DRS sub modules are defined. They are called from the DRS in the following manner. \\

\begin{pythonbox}
from SpirouDRS import spirouX
spirouX.MyFunction(arg1, arg2, kwarg1='value')   # Alias to internal function
spirouX.spirouY.my_function(arg1, arg2, kwarg1='value')   # internal function
\end{pythonbox}

\noindent Described in each section are the functions used in the recipes (defined in each sub-modules \path{__init__.py} file) and as such are in CamelCase (capitalized). The internal links are provided to locate the code for each function but only the CamelCased aliases should be used in recipes. 

Functions used with in these sub-modules that are used internally but not called from recipes are not defined here but all functions are described as with those below and can be read in any interactive python or ipython terminal in the following way: \\

\noindent\begin{minipage}{\textwidth}
\begin{pythonbox}
# in python
from SpirouDRS import spirouConfig
print(spirouConfig.CheckConfig.__doc__)
\end{pythonbox}
\vspace{-0.25cm}
\begin{cmdboxprint}
Check whether we have certain keys in dictionary
raises a Config Error if keys are not in params
:param p: parameter dictionary, ParamDict containing constants
    Must contain at least:
        the keys defined in "keys" (else ConfigError raised)
:param keys: string or list of strings containing the keys to look for
:return None:
\end{cmdboxprint}
\end{minipage}

\vspace{0.5cm}

\noindent\begin{minipage}{\textwidth}
\begin{pythonbox}
# in ipython
from SpirouDRS import spirouConfig
spirouConfig.CheckConfig?
\end{pythonbox}
\vspace{-0.25cm}
\begin{cmdboxprintspecial}
@rSignature:@r spirouConfig@b.@bCheckConfig@b(@bparams@b,@b keys@b)@b
@rDocstring:@r
Check whether we have certain keys in dictionary
raises a Config Error if keys are not in params
:param p: parameter dictionary, ParamDict containing constants
    Must contain at least:
        the keys defined in "keys" (else ConfigError raised)
:param keys: string or list of strings containing the keys to look for
:return None:
@rFile:@r      /drs/INTROOT/SpirouDRS/spirouConfig/spirouConfig.py
@rType:@r      function
\end{cmdboxprintspecial}
\end{minipage}

%%%%%%%%%%%%%%%%%%%%%%%%%%%%%%%%%%%%%%%%%%%%%%%%%%%%%%%%
%%
%%    Modules imported from tex files
%%
%%%%%%%%%%%%%%%%%%%%%%%%%%%%%%%%%%%%%%%%%%%%%%%%%%%%%%%%


%%%%%%%%%%%%%%%%%%%%%%%%%%%%%%%%%%%%%%%%%%%%%%%%%%%%%%%%
%%
\clearpage
\newpage
\noindent\begin{minipage}{\textwidth}
\section{The spirouBACK module}
\label{ch:the_module:spirouBACK}
%%
%%%%%%%%%%%%%%%%%%%%%%%%%%%%%%%%%%%%%%%%%%%%%%%%%%%%%%%%

%----------------------------------------------------------------------------------------

\subsection{BoxSmoothedMinMax}
Defined in \spirouBACK\path{.measure_box_min_max}
\begin{pythonbox}
from SpirouDRS import spirouBACK
spirouBACK.BoxSmoothedMinMax(y, size)
spirouBACK.spirouBACK.measure_box_min_max(y, size)
\end{pythonbox}
\begin{pythondocstring}
Measure the minimum and maximum pixel value for each pixel using a box which
surrounds that pixel by:  pixel-size to pixel+size.

Edge pixels (0-->size and (len(y)-size)-->len(y) are
set to the values for pixel=size and pixel=(len(y)-size)

:param y: numpy array (1D), the image
:param size: int, the half size of the box to use (half height)
             so box is defined from  pixel-size to pixel+size

:return min_image: numpy array (1D length = len(y)), the values
                   for minimum pixel defined by a box of pixel-size to
                   pixel+size for all columns
:return max_image: numpy array (1D length = len(y)), the values
                   for maximum pixel defined by a box of pixel-size to
                   pixel+size for all columns
\end{pythondocstring}
\end{minipage}


%----------------------------------------------------------------------------------------
\noindent\begin{minipage}{\textwidth}
\subsection{MeasureBackgroundFF}

Defined in \spirouBACK\path{.measure_background_flatfield}

\begin{pythonbox}
from SpirouDRS import spirouBACK
spirouBACK.MeasureBackgroundFF(p, image)
spirouBACK.spirouBACK.measure_background_flatfield(p, image)
\end{pythonbox}

\begin{pythondocstring}
Measures the background of a flat field image - currently does not work
as need an interpolation function (see code)

:param p: parameter dictionary, ParamDict containing constants

        Must contain at least:
            IC_BKGR_WINDOW: int, Half-size of window for background
                            measurements
            GAIN: float, the gain of the image (from HEADER)
            SIGDET: float, the read noise of the image (from HEADER)
            log_opt: string, log option, normally the program name

:param image: numpy array (2D), the image to measure the background of

:return background: numpy array (2D), the background image (currently all
                    zeros) as background not implemented
:return xc: numpy array (1D), the box centres (x positions) used to create
            the background image
:return yc: numpy array (1D), the box centres (y positions) used to create
            the background image
:return minlevel: numpy array (2D), the 2 * size -th minimum pixel value
                  of each box for each pixel in the image
\end{pythondocstring}
\end{minipage}

%----------------------------------------------------------------------------------------
\noindent\begin{minipage}{\textwidth}
\subsection{MeasureMinMax}

Defined in \spirouBACK\path{measure_box_min_max}

\begin{pythonbox}
from SpirouDRS import spirouBACK
spirouBACK.MeasureMinMax(y, size)
spirouBACK.spirouBACK.measure_box_min_max(y, size)
\end{pythonbox}

\begin{pythondocstring}
Measure the minimum and maximum pixel value for each pixel using a box which
surrounds that pixel by:  pixel-size to pixel+size.

Edge pixels (0-->size and (len(y)-size)-->len(y) are
set to the values for pixel=size and pixel=(len(y)-size)

:param y: numpy array (1D), the image
:param size: int, the half size of the box to use (half height)
             so box is defined from  pixel-size to pixel+size

:return min_image: numpy array (1D length = len(y)), the values
                   for minimum pixel defined by a box of pixel-size to
                   pixel+size for all columns
:return max_image: numpy array (1D length = len(y)), the values
                   for maximum pixel defined by a box of pixel-size to
                   pixel+size for all columns
\end{pythondocstring}
\end{minipage}

%----------------------------------------------------------------------------------------
\noindent\begin{minipage}{\textwidth}
\subsection{MeasureMinMaxSignal}

Defined in \spirouBACK\path{measure_min_max}

\begin{pythonbox}
from SpirouDRS import spirouBACK
spirouBACK.MeasureMinMaxSignal(pp, y)
spirouBACK.spirouBACK.measure_min_max(pp, y)
\end{pythonbox}

\begin{pythondocstring}
Measure the minimum, maximum peak to peak values in y, the third biggest
pixel in y and the peak-to-peak difference between the minimum and
maximum values in y

:param pp: parameter dictionary, ParamDict containing constants
            Must contain at least:
                IC_LOCNBPIX: int, Half spacing between orders

:param y: numpy array (1D), the central column pixel values

:return miny: numpy array (1D length = len(y)), the values
              for minimum pixel defined by a box of pixel-size to
              pixel+size for all columns
:return maxy: numpy array (1D length = len(y)), the values
              for maximum pixel defined by a box of pixel-size to
              pixel+size for all columns
:return max_signal: float, the pixel value of the third biggest value
                    in y
:return diff_maxmin: float, the difference between maxy and miny
\end{pythondocstring}
\end{minipage}


%----------------------------------------------------------------------------------------
\noindent\begin{minipage}{\textwidth}
\subsection{MeasureBkgrdGetCentPixs}

Defined in \spirouBACK\path{measure_background_and_get_central_pixels}

\begin{pythonbox}
from SpirouDRS import spirouBACK
spirouBACK.MeasureBkgrdGetCentPixs(pp, loc, image)
spirouBACK.spirouBACK.measure_background_and_get_central_pixels(pp, loc, image)
\end{pythonbox}

\begin{pythondocstring}
Takes the image and measure the background

:param pp: parameter dictionary, ParamDict containing constants
        Must contain at least:
            IC_OFFSET: int, row number of image to start processing at
            IC_CENT_COL: int, Definition of the central column
            IC_MIN_AMPLITUDE: int, Minimum amplitude to accept (in e-)
            IC_LOCSEUIL: float, Normalised amplitude threshold to accept
                         pixels for background calculation
            log_opt: string, log option, normally the program name
            DRS_DEBUG: int, Whether to run in debug mode
                            0: no debug
                            1: basic debugging on errors
                            2: recipes specific (plots and some code runs)
            DRS_PLOT: bool, Whether to plot (True to plot)

:param loc: parameter dictionary, ParamDict containing data

:param image: numpy array (2D), the image

:return ycc: the normalised values the central pixels

:return loc: parameter dictionary, the updated parameter dictionary
        Adds/updates the following:
            ycc: numpy array (1D), normalized central column of pixels
            mean_backgrd: float, 100 times the mean of the good background
                          pixels
            max_signal: float, the maximum value of the central column of
                        pixels
\end{pythondocstring}
\end{minipage}
%----------------------------------------------------------------------------------------

\input{Chapters/modules/spirouDB}


%%%%%%%%%%%%%%%%%%%%%%%%%%%%%%%%%%%%%%%%%%%%%%%%%%%%%%%%
%%
\clearpage
\newpage
\begin{minipage}{\textwidth}
\section{The  module}
\label{ch:the_module:}
%%
%%%%%%%%%%%%%%%%%%%%%%%%%%%%%%%%%%%%%%%%%%%%%%%%%%%%%%%%

%----------------------------------------------------------------------------------------
\subsection{ConfigError}

Defined in \spirou \path{.}

\begin{pythonbox}
from SpirouDRS import spirou
spirouConfig.ConfigError

\end{pythonbox}

\begin{pythondocstring}

\end{pythondocstring}
\end{minipage}

%----------------------------------------------------------------------------------------
\begin{minipage}{\textwidth}
\subsection{CheckCparams}

Defined in \spirou \path{.}

\begin{pythonbox}
from SpirouDRS import spirou
spirouConfig.CheckCparams

\end{pythonbox}

\begin{pythondocstring}

\end{pythondocstring}
\end{minipage}

%----------------------------------------------------------------------------------------
\begin{minipage}{\textwidth}
\subsection{CheckConfig}

Defined in \spirou \path{.}

\begin{pythonbox}
from SpirouDRS import spirou
spirouConfig.CheckConfig

\end{pythonbox}

\begin{pythondocstring}

\end{pythondocstring}
\end{minipage}

%----------------------------------------------------------------------------------------
\begin{minipage}{\textwidth}
\subsection{ExtractDictParams}

Defined in \spirou \path{.}

\begin{pythonbox}
from SpirouDRS import spirou
spirouConfig.ExtractDictParams

\end{pythonbox}

\begin{pythondocstring}

\end{pythondocstring}
\end{minipage}

%----------------------------------------------------------------------------------------
\begin{minipage}{\textwidth}
\subsection{GetKeywordArguments}

Defined in \spirou \path{.}

\begin{pythonbox}
from SpirouDRS import spirou
spirouConfig.GetKeywordArguments

\end{pythonbox}

\begin{pythondocstring}

\end{pythondocstring}
\end{minipage}

%----------------------------------------------------------------------------------------
\begin{minipage}{\textwidth}
\subsection{GetKwValues}

Defined in \spirou \path{.}

\begin{pythonbox}
from SpirouDRS import spirou
spirouConfig.GetKwValues

\end{pythonbox}

\begin{pythondocstring}

\end{pythondocstring}
\end{minipage}

%----------------------------------------------------------------------------------------
\begin{minipage}{\textwidth}
\subsection{GetAbsFolderPath}

Defined in \spirou \path{.}

\begin{pythonbox}
from SpirouDRS import spirou
spirouConfig.GetAbsFolderPath

\end{pythonbox}

\begin{pythondocstring}

\end{pythondocstring}
\end{minipage}

%----------------------------------------------------------------------------------------
\begin{minipage}{\textwidth}
\subsection{GetDefaultConfigFile}

Defined in \spirou \path{.}

\begin{pythonbox}
from SpirouDRS import spirou
spirouConfig.GetDefaultConfigFile

\end{pythonbox}

\begin{pythondocstring}

\end{pythondocstring}
\end{minipage}

%----------------------------------------------------------------------------------------
\begin{minipage}{\textwidth}
\subsection{LoadConfigFromFile}

Defined in \spirou \path{.}

\begin{pythonbox}
from SpirouDRS import spirou
spirouConfig.LoadConfigFromFile

\end{pythonbox}

\begin{pythondocstring}

\end{pythondocstring}
\end{minipage}

%----------------------------------------------------------------------------------------
\begin{minipage}{\textwidth}
\subsection{ParamDict}

Defined in \spirou \path{.}

\begin{pythonbox}
from SpirouDRS import spirou
spirouConfig.ParamDict

\end{pythonbox}

\begin{pythondocstring}

\end{pythondocstring}
\end{minipage}

%----------------------------------------------------------------------------------------
\begin{minipage}{\textwidth}
\subsection{ReadConfigFile}

Defined in \spirou \path{.}

\begin{pythonbox}
from SpirouDRS import spirou
spirouConfig.ReadConfigFile

\end{pythonbox}

\begin{pythondocstring}

\end{pythondocstring}
\end{minipage}

%----------------------------------------------------------------------------------------


%%%%%%%%%%%%%%%%%%%%%%%%%%%%%%%%%%%%%%%%%%%%%%%%%%%%%%%%
%%
\clearpage
\newpage
\noindent\begin{minipage}{\textwidth}
\section{The spirouCore module}
\label{ch:the_module:spirouCore}
%%
%%%%%%%%%%%%%%%%%%%%%%%%%%%%%%%%%%%%%%%%%%%%%%%%%%%%%%%%

%----------------------------------------------------------------------------------------
\subsection{wlog}
\label{ch:the_module:spirouCore:logger}

Defined in \spirouCore\path{.spirouLog.logger}, also aliased in code to `WLOG'.

See Section \ref{ch:rules:drs_specific:logger} for usage details.

\begin{pythonbox}
from SpirouDRS import spirouCore
spirouCore.wlog(key='', option='', message='', printonly=False, logonly=False)
spirouCore.spirouLog.logger(key='', option='', message='', printonly=False, logonly=False)
\end{pythonbox}

\begin{pythondocstring}
Parses a key (error/warning/info/graph), an option and a message to the
stdout and the log file.

keys are controlled by "spirouConfig.Constants.LOG_TRIG_KEYS()"
printing to screen is controlled by "PRINT_LEVEL" constant (config.py)
printing to log file is controlled by "LOG_LEVEL" constant (config.py)
based on the levels described in "spirouConfig.Constants.WRITE_LEVEL"

:param key: string, either "error" or "warning" or "info" or graph, this
            gives a character code in output
:param option: string, option code
:param message: string or list of strings, message to display or messages
                to display (1 line for each message in list)
:param printonly: bool, print only do not save to log (default = False)
:param logonly: bool, log only do not save to log (default = False)

output to stdout/log is as follows:

    HH:MM:SS.S - CODE |option|message

time is output in UTC to nearest .1 seconds

:return None:
\end{pythondocstring}
\end{minipage}

%----------------------------------------------------------------------------------------
\noindent\begin{minipage}{\textwidth}
\subsection{WarnLog}

Defined in \spirouCore\path{.spirouLog.warninglogger}

\begin{pythonbox}
from SpirouDRS import spirouCore
spirouCore.warnlog(w, funcname=None)
spirouCore.spirouLog.warninglogger(w, funcname=None)
\end{pythonbox}

\begin{pythondocstring}
Warning logger - takes "w" - a list of caught warnings and pipes them on
to the log functions.
If "funcname" is not None then "funcname" is printed with the line reference
(intended to be used to identify the code/function/module warning was
generated in)

to catch warnings use the following:

>>> import warnings
>>> with warnings.catch_warnings(record=True) as w:
>>>     code_to_generate_warnings()
>>> warninglogger(w, 'some name for logging')

:param w: list of warnings, the list of warnings from
           warnings.catch_warnings
:param funcname: string or None, if string then also pipes "funcname" to the
                 warning message (intended to be used to identify the code/
                 function/module warning was generated in)
:return:
\end{pythondocstring}
\end{minipage}

%----------------------------------------------------------------------------------------
\noindent\begin{minipage}{\textwidth}
\subsection{GaussFunction}

Defined in \spirouCore\path{.spirouMath.gauss_function}

\begin{pythonbox}
from SpirouDRS import spirouCore
spirouCore.GaussFunction(x, a, x0, sigma, dc)
spirouCore.spirouMath.gauss_function(x, a, x0, sigma, dc)
\end{pythonbox}

\begin{pythondocstring}
A standard 1D Gaussian function (for fitting against)]=

:param x: numpy array (1D), the x data points
:param a: float, the amplitude
:param x0: float, the mean of the Gaussian
:param sigma: float, the standard deviation (FWHM) of the Gaussian
:param dc: float, the constant level below the Gaussian

:return gauss: numpy array (1D), size = len(x), the output Gaussian
\end{pythondocstring}
\end{minipage}

%----------------------------------------------------------------------------------------
\noindent\begin{minipage}{\textwidth}
\subsection{GetTimeNowUnix}

Defined in \spirouCore\path{.spirouMath.get_time_now_unix}

\begin{pythonbox}
from SpirouDRS import spirouCore
spirouCore.GetTimeNowUnix(zone='UTC')
spirouCore.spirouMath.get_time_now_unix(zone='UTC')
\end{pythonbox}

\begin{pythondocstring}
Get the unix_time now.

Default is to return unix_time in UTC/GMT time

:param zone: string, if UTC displays the time in UTC else displays local
             time

:return unix_time: float, the unix_time
\end{pythondocstring}
\end{minipage}

%----------------------------------------------------------------------------------------
\noindent\begin{minipage}{\textwidth}
\subsection{GetTimeNowString}

Defined in \spirouCore\path{.spirouMath.get_time_now_string}

\begin{pythonbox}
from SpirouDRS import spirouCore
spirouCore.GetTimeNowString(fmt=TIME_FMT, zone='UTC')
spirouCore.spirouMath.get_time_now_string(fmt=TIME_FMT, zone='UTC')
\end{pythonbox}

\begin{pythondocstring}
Get the time now (in string format = "fmt")

Default is to return string time in UTC/GMT time

    Commonly used format codes:

    %Y  Year with century as a decimal number.
    %m  Month as a decimal number [01,12].
    %d  Day of the month as a decimal number [01,31].
    %H  Hour (24-hour clock) as a decimal number [00,23].
    %M  Minute as a decimal number [00,59].
    %S  Second as a decimal number [00,61].
    %z  Time zone offset from UTC.
    %a  Locale's abbreviated weekday name.
    %A  Locale's full weekday name.
    %b  Locale's abbreviated month name.
    %B  Locale's full month name.
    %c  Locale's appropriate date and time representation.
    %I  Hour (12-hour clock) as a decimal number [01,12].
    %p  Locale's equivalent of either AM or PM.


:param fmt: string, the format code for the returned time
:param zone: string, if UTC displays the time in UTC else displays local
             time

:return stringtime: string, the time in a string in format = "fmt"
\end{pythondocstring}
\end{minipage}


%----------------------------------------------------------------------------------------
\noindent\begin{minipage}{\textwidth}
\subsection{PrintLog}

Defined in \spirouCore\path{.spirouMath.printlog}

\begin{pythonbox}
from SpirouDRS import spirouCore
spirouCore.PrintLog(message, key='all')
spirouCore.spirouCore.printlog(message, key='all')
\end{pythonbox}

\begin{pythondocstring}
print message to stdout (if level is correct - set by PRINT_LEVEL)
is coloured unless spirouConfig.Constants.COLOURED_LOG() is False

:param message: string, the formatted log line to write to stdout
:param key: string, either "error" or "warning" or "info" or "graph" or
            "all", this gives a character code in output

:return None:
\end{pythondocstring}
\end{minipage}


%----------------------------------------------------------------------------------------
\noindent\begin{minipage}{\textwidth}
\subsection{PrintColours}

Defined in \spirouCore\path{.spirouMath.printlog}

\begin{pythonbox}
from SpirouDRS import spirouCore
spirouCore.PrintColour(key='all', func_name=None)
spirouCore.spirouCore.printcolour(key='all', func_name=None)
\end{pythonbox}

\begin{pythondocstring}
Get the print colour (start and end) based on "key".
This should be used as follows:
    >> c1, c2 = printcolour(key='all')
    >> print(c1 + message + c2)

:param key: string, either "error" or "warning" or "info" or "graph" or
            "all", this gives a character code in output
:param func_name: string or None, if not None then defines the function to
                  report in the error
:return colour1: string or None, if key is found and we are using coloured
                 log returns the starting colour, if not returns empty
                 string if key is not accepted does not print
:return colour2: string or None, if key is found and we are using coloured
                 log returns the ending colour, if not returns empty
                 string if key is not accepted does not print
\end{pythondocstring}
\end{minipage}



%----------------------------------------------------------------------------------------
\noindent\begin{minipage}{\textwidth}
\subsection{Unix2stringTime}

Defined in \spirouCore\path{.spirouMath.unixtime2stringtime}

\begin{pythonbox}
from SpirouDRS import spirouCore
spirouCore.Unix2stringTime(ts, fmt=DATE_FMT, zone='UTC')
spirouCore.spirouMath.unixtime2stringtime(ts, fmt=DATE_FMT, zone='UTC')
\end{pythonbox}

\begin{pythondocstring}
Convert a unix time (seconds since  1970-01-01 00:00:00 GMT) into a
string in format "fmt". Currently supported time-zones are UTC and local
(i.e. your current time zone).

Default is to return string time in UTC/GMT time

Commonly used format codes:

    %Y  Year with century as a decimal number.
    %m  Month as a decimal number [01,12].
    %d  Day of the month as a decimal number [01,31].
    %H  Hour (24-hour clock) as a decimal number [00,23].
    %M  Minute as a decimal number [00,59].
    %S  Second as a decimal number [00,61].
    %z  Time zone offset from UTC.
    %a  Locale's abbreviated weekday name.
    %A  Locale's full weekday name.
    %b  Locale's abbreviated month name.
    %B  Locale's full month name.
    %c  Locale's appropriate date and time representation.
    %I  Hour (12-hour clock) as a decimal number [01,12].
    %p  Locale's equivalent of either AM or PM.

:param ts: float or int, the unix time (seconds since 1970-01-01 00:00:00
           GMT)
:param fmt: string, the format of the string to convert
:param zone: string, the time zone for the input string
                      (currently supported =  "UTC" or "local")

:return stringtime: string, the time in format "fmt"
\end{pythondocstring}
\end{minipage}

%----------------------------------------------------------------------------------------
\noindent\begin{minipage}{\textwidth}
\subsection{String2unixTime}

Defined in \spirouCore\path{.spirouMath.stringtime2unixtime}

\begin{pythonbox}
from SpirouDRS import spirouCore
spirouCore.String2unixTime(string, fmt=DATE_FMT, zone='UTC')
spirouCore.spirouMath.stringtime2unixtime(string, fmt=DATE_FMT, zone='UTC')
\end{pythonbox}

\begin{pythondocstring}
Convert a string in format "fmt" into a float Unix time (seconds since
1970-01-01 00:00:00 GMT). Currently supported time-zones are UTC and local
(i.e. your current time zone).

Default is to assume string is in UTC/GMT time

Commonly used format codes:

    %Y  Year with century as a decimal number.
    %m  Month as a decimal number [01,12].
    %d  Day of the month as a decimal number [01,31].
    %H  Hour (24-hour clock) as a decimal number [00,23].
    %M  Minute as a decimal number [00,59].
    %S  Second as a decimal number [00,61].
    %z  Time zone offset from UTC.
    %a  Locale's abbreviated weekday name.
    %A  Locale's full weekday name.
    %b  Locale's abbreviated month name.
    %B  Locale's full month name.
    %c  Locale's appropriate date and time representation.
    %I  Hour (12-hour clock) as a decimal number [01,12].
    %p  Locale's equivalent of either AM or PM.

:param string: string, the time string to convert
:param fmt: string, the format of the string to convert
:param zone: string, the time zone for the input string
                      (currently supported =  "UTC" or "local")

:return unix_time: float, Unix time (seconds since 1970-01-01 00:00:00 GMT)
\end{pythondocstring}
\end{minipage}

%----------------------------------------------------------------------------------------
\noindent\begin{minipage}{\textwidth}
\subsection{sPlt}

Defined in \spirouCore\path{.spirouPlot} (alias to the plotting module).

\begin{pythonbox}
from SpirouDRS import spirouCore
spirouCore.sPlt
spirouCore.spirouPlot
\end{pythonbox}

\begin{pythondocstring}
Spirou Plotting functions available:

start_interactive_session(interactive=False)
    Start interactive plot session, if required and if
    spirouConfig.Constants.INTERACITVE_PLOTS_ENABLED() is True

    :param interactive: bool, if True start interactive session

    :return None:

end_interactive_session(interactive=False)
    End interactive plot session, if required and if
    spirouConfig.Constants.INTERACITVE_PLOTS_ENABLED() is True

    :param interactive: bool, if True end interactive session

    :return None:


define_figure
    Define a figure number (mostly for use in interactive mode)

    :param num: int, a figure number

    :return figure: plt.figure instance

closeall()
    Close all matplotlib plots currently open

    :return None:

And all plotting functions from specific recipes.
\end{pythondocstring}
\end{minipage}

%----------------------------------------------------------------------------------------

\noindent\begin{minipage}{\textwidth}
\subsection{Getll}

Defined in \spirouCore\path{.spirouMath.get_ll_from_coefficients}

\begin{pythonbox}
from SpirouDRS import spirouCore
spirouCore.Getll(params, nx, nbo)
spirouCore.spirouMath.get_ll_from_coefficients(params, nx, nbo)
\end{pythonbox}

\begin{pythondocstring}
Use the coefficient matrix "params" to construct fit values for each order
(dimension 0 of coefficient matrix) for values of x from 0 to nx
(interger steps)

:param params: numpy array (2D), the coefficient matrix
               size = (number of orders x number of fit coefficients)

:param nx: int, the number of values and the maximum value of x to use
           the coefficients for, where x is such that

            yfit = p[0]*x**(N-1) + p[1]*x**(N-2) + ... + p[N-2]*x + p[N-1]

            N = number of fit coefficients
            and p is the coefficients for one order
            (i.e. params = [ p_1, p_2, p_3, p_4, p_5, ... p_nbo]

:param nbo: int, the number of orders to use

:return ll: numpy array (2D): the yfit values for each order
            (i.e. ll = [yfit_1, yfit_2, yfit_3, ..., yfit_nbo] )
\end{pythondocstring}
\end{minipage}

%----------------------------------------------------------------------------------------
\noindent\begin{minipage}{\textwidth}
\subsection{Getdll}

Defined in \spirouTHORCA\path{.spirouTHORCA.get_dll_from_coefficients}

\begin{pythonbox}
from SpirouDRS import spirouCore
spirouCore.Getdll(params, nx, nbo)
spirouCore.spirouMath.get_dll_from_coefficients(params, nx, nbo)
\end{pythonbox}

\begin{pythondocstring}
Derivative of the coefficients, using the coefficient matrix "params"
to construct the derivative of the fit values for each order
(dimension 0 of coefficient matrix) for values of x from 0 to nx
(interger steps)

:param params: numpy array (2D), the coefficient matrix
               size = (number of orders x number of fit coefficients)

:param nx: int, the number of values and the maximum value of x to use
           the coefficients for, where x is such that

            yfit = p[0]*x**(N-1) + p[1]*x**(N-2) + ... + p[N-2]*x + p[N-1]

            dyfit = p[0]*(N-1)*x**(N-2) + p[1]*(N-2)*x**(N-3) + ... +
                    p[N-3]*x + p[N-2]

            N = number of fit coefficients
            and p is the coefficients for one order
            (i.e. params = [ p_1, p_2, p_3, p_4, p_5, ... p_nbo]

:param nbo: int, the number of orders to use

:return ll: numpy array (2D): the yfit values for each order
            (i.e. ll = [dyfit_1, dyfit_2, dyfit_3, ..., dyfit_nbo] )
\end{pythondocstring}
\end{minipage}


%%%%%%%%%%%%%%%%%%%%%%%%%%%%%%%%%%%%%%%%%%%%%%%%%%%%%%%%
%%
\clearpage
\newpage
\begin{minipage}{\textwidth}
\section{The spirouEXTOR module}
\label{ch:the_module:spirouEXTOR}
%%
%%%%%%%%%%%%%%%%%%%%%%%%%%%%%%%%%%%%%%%%%%%%%%%%%%%%%%%%

%----------------------------------------------------------------------------------------
\subsection{Extraction}

Defined in \spirouEXTOR\path{.spirouEXTOR.extract_wrapper}

\begin{pythonbox}
from SpirouDRS import spirou
spirouEXTOR.Extraction(image, pos, sig, **kwargs)
spirouEXTOR.spirouEXTOR.extract_wrapper(image, pos, sig, **kwargs)
\end{pythonbox}

\begin{pythondocstring}
Extraction wrapper - takes in image, pos, sig and kwargs and decides
which extraction process to use.

:param image: numpy array (2D), the image
:param pos: numpy array (1D), the position fit coefficients
            size = number of coefficients for fit
:param sig: numpy array (1D), the width fit coefficients
            size = number of coefficients for fit
:param kwargs: additional keyword arguments

currently accepted keyword arguments are:

    extopt:         int, Extraction option in tilt file:
                     if 0 extraction by summation over constant range
                     if 1 extraction by summation over constant sigma
                        (not currently available)
                     if 2 Horne extraction without cosmic elimination
                        (not currently available)
                     if 3 Horne extraction with cosmic elimination
                        (not currently available)

    nbsig:          float,  distance away from center to extract out to +/-
                    defaults to p['nbsig'] from constants_SPIROU.py

    gain:           float, gain of the image
                    defaults to p['gain'] from fitsfilename HEADER

    sigdet:         float, the sigdet of the image
                    defaults to p['sigdet'] from fitsfilename HEADER

    range1:         float, Half-zone extraction width left side
                    (formally plage1)
                    defaults to p['ic_ext_range1'] from fiber parameters in
                    constatns_SPIROU.txt

    range2:         float, Half-zone extraction width left side
                    (formally plage2)
                    defaults to p['ic_ext_range2'] from fiber parameters in
                    constatns_SPIROU.txt

    tilt:           numpy array (1D), the tilt for this order, if defined
                    uses tilt, if not defined does not
\end{pythondocstring}
\end{minipage}

\begin{minipage}{\textwidth}
\begin{pythondocstring}
Extraction wrapper (continued)

currently accepted keyword arguments are: (continued)

    use_weight:    bool, if True use weighted extraction, if False or not
                    defined does not use weighted extraction

    order_profile:  numpy array (2D), the image with fit superposed on top,
                    required for tilt and or weighted fit

    mode:           if use_weight and tilt is not None then
                    if mode = 'old'  will use old code (use this if
                    exception generated)
                    extract_tilt_weight_order_old() is run

                    else mode = 'new' and
                    extract_tilt_weight_order() is run

:return spe: numpy array (1D), the extracted pixel values,
             size = image.shape[1] (along the order direction)
:return nbcos: int, zero in this case
\end{pythondocstring}
\end{minipage}

%----------------------------------------------------------------------------------------
\begin{minipage}{\textwidth}
\subsection{ExtractABOrderOffset}

Defined in \spirouEXTOR\path{.spirouEXTOR.extract_AB_order}

\begin{pythonbox}
from SpirouDRS import spirou
spirouEXTOR.ExtractABOrderOffset(pp, loc, image, rnum)
spirouEXTOR.spirouEXTOR.extract_AB_order(pp, loc, image, rnum)
\end{pythonbox}

\begin{pythondocstring}
Perform the extraction on the AB fibers separately using the summation
over constant range

:param p: parameter dictionary, ParamDict containing constants
    Must contain at least:
            IC_CENT_COL: int, the column number (x-axis) of the central
                         column
            IC_FACDEC: float, the offset multiplicative factor for width
            IC_EXTOPT: int, the extraction option
            gain: float, the gain of the image
            IC_EXTNBSIG: float, distance away from center to extract
                         out to +/- (in rows or y-axis direction)

:param loc: parameter dictionary, ParamDict containing data
        Must contain at least:
            ass: numpy array (2D), the fit coefficients array for
                  the widths fit
                  shape = (number of orders x number of fit coefficients)
            acc: numpy array (2D), the fit coefficients array for
                  the centers fit
                  shape = (number of orders x number of fit coefficients)


:param image: numpy array (2D), the image
:param rnum: int, the order number for this iteration

:return loc: parameter dictionary, the updated parameter dictionary
        Adds/updates the following:
            offset: numpy array (1D), the center values with the
                    offset in 'IC_CENT_COL' added
            cent1: numpy array (2D), the extraction for A, updated is
                   the order "rnum"
            nbcos: int, 0 (constant)
            cent2: numpy array (2D), the extraction for B, updated is
                   the order "rnum"
\end{pythondocstring}
\end{minipage}

%----------------------------------------------------------------------------------------
\begin{minipage}{\textwidth}
\subsection{ExtractOrder}

Defined in \spirouEXTOR\path{.spirouEXTOR.extract_order}

\begin{pythonbox}
from SpirouDRS import spirou
spirouEXTOR.ExtractOrder(pp, loc, image, rnum, **kwargs)
spirouEXTOR.spirouEXTOR.extract_order(pp, loc, image, rnum, **kwargs)
\end{pythonbox}

\begin{pythondocstring}
Extract order without tilt or weight using spirouEXTOR.extract_wrapper()

:param p: parameter dictionary, ParamDict containing constants
    Must contain at least:
            IC_EXTOPT: int, the extraction option
            IC_EXT_RANGE: float, the upper and lower edge of the order
                          in rows (y-axis) - half-zone width
            gain: float, the gain of the image

:param loc: parameter dictionary, ParamDict containing data
        Must contain at least:
            acc: numpy array (2D), the fit coefficients array for
                  the centers fit
                  shape = (number of orders x number of fit coefficients)
            ass: numpy array (2D), the fit coefficients array for
                  the widths fit
                  shape = (number of orders x number of fit coefficients)

:param image: numpy array (2D), the image
:param rnum: int, the order number for this iteration
:param kwargs: additional keywords to pass to the extraction wrapper

        - allowed keywords are:

        range1  (defaults to "IC_EXT_RANGE")
        range2  (defaults to "IC_EXT_RANGE")
        gain    (defaults to "GAIN")

:return cent: numpy array (1D), the extracted pixel values,
             size = image.shape[1] (along the order direction)
:return cpt: int, zero in this case
\end{pythondocstring}
\end{minipage}

%----------------------------------------------------------------------------------------

\begin{minipage}{\textwidth}
\subsection{ExtractTiltOrder}

Defined in \spirouEXTOR\path{.spirouEXTOR.extract_tilt_order}

\begin{pythonbox}
from SpirouDRS import spirou
spirouEXTOR.ExtractTiltOrder(pp, loc, image, rnum, **kwargs)
spirouEXTOR.spirouEXTOR.extract_tilt_order(pp, loc, image, rnum, **kwargs)
\end{pythonbox}

\begin{pythondocstring}
Extract order with tilt but without weight using
spirouEXTOR.extract_wrapper()

:param p: parameter dictionary, ParamDict containing constants
    Must contain at least:
            IC_EXT_RANGE: float, the upper and lower edge of the order
                          in rows (y-axis) - half-zone width
            gain: float, the gain of the image

:param loc: parameter dictionary, ParamDict containing data
        Must contain at least:
            acc: numpy array (2D), the fit coefficients array for
                  the centers fit
                  shape = (number of orders x number of fit coefficients)
            ass: numpy array (2D), the fit coefficients array for
                  the widths fit
                  shape = (number of orders x number of fit coefficients)
            tilt: numpy array (1D), the tilt angle of each order

:param image: numpy array (2D), the image
:param rnum: int, the order number for this iteration
:param kwargs: additional keywords to pass to the extraction wrapper

        - allowed keywords are:

        range1  (defaults to "IC_EXT_RANGE")
        range2  (defaults to "IC_EXT_RANGE")
        gain    (defaults to "GAIN")

:return cent: numpy array (1D), the extracted pixel values,
             size = image.shape[1] (along the order direction)
:return cpt: int, zero in this case
\end{pythondocstring}
\end{minipage}

%----------------------------------------------------------------------------------------

\begin{minipage}{\textwidth}
\subsection{ExtractTiltWeightOrder}

Defined in \spirouEXTOR\path{.extract_tilt_weight_order}

\begin{pythonbox}
from SpirouDRS import spirou
spirouEXTOR.ExtractTiltWeightOrder(pp, loc, image, orderp, rnum, **kwargs)
spirouEXTOR.spirouEXTOR.extract_tilt_weight_order(pp, loc, image, orderp, rnum, **kwargs)
\end{pythonbox}

\begin{pythondocstring}
Extract order with tilt and weight using
spirouEXTOR.extract_wrapper() with mode=1
(extract_tilt_weight_order_old() is run)

:param p: parameter dictionary, ParamDict containing constants
    Must contain at least:
            IC_EXT_RANGE: float, the upper and lower edge of the order
                          in rows (y-axis) - half-zone width
            gain: float, the gain of the image
            sigdet: float, the read noise of the image

:param loc: parameter dictionary, ParamDict containing data
        Must contain at least:
            acc: numpy array (2D), the fit coefficients array for
                  the centers fit
                  shape = (number of orders x number of fit coefficients)
            ass: numpy array (2D), the fit coefficients array for
                  the widths fit
                  shape = (number of orders x number of fit coefficients)
            tilt: numpy array (1D), the tilt angle of each order

:param image: numpy array (2D), the image
:param orderp: numpy array (2D), the order profile image
:param rnum: int, the order number for this iteration
:param kwargs: additional keywords to pass to the extraction wrapper

        - allowed keywords are:

        range1  (defaults to "IC_EXT_RANGE")
        range2  (defaults to "IC_EXT_RANGE")
        gain    (defaults to "GAIN")
        sigdet  (defaults to "SIGDET")

:return cent: numpy array (1D), the extracted pixel values,
             size = image.shape[1] (along the order direction)
:return cpt: int, zero in this case
\end{pythondocstring}
\end{minipage}

%----------------------------------------------------------------------------------------

\begin{minipage}{\textwidth}
\subsection{ExtractTiltWeightOrder2}

Defined in \spirouEXTOR\path{.}

\begin{pythonbox}
from SpirouDRS import spirou
spirouEXTOR.ExtractTiltWeightOrder2(pp, loc, image, orderp, rnum, **kwargs)
spirouEXTOR.spirouEXTOR.extract_tilt_weight_order2(pp, loc, image, orderp, rnum, **kwargs)
\end{pythonbox}

\begin{pythondocstring}
Extract order with tilt and weight using
spirouEXTOR.extract_wrapper() with mode=2
(extract_tilt_weight_order() is run)

:param p: parameter dictionary, ParamDict containing constants
    Must contain at least:
            IC_EXT_RANGE1: float, the upper edge of the order in rows
                           (y-axis) - half-zone width (lower)
            IC_EXT_RANGE2: float, the lower edge of the order in rows
                           (y-axis) - half-zone width (upper)
            gain: float, the gain of the image
            sigdet: float, the read noise of the image

:param loc: parameter dictionary, ParamDict containing data
        Must contain at least:
            acc: numpy array (2D), the fit coefficients array for
                  the centers fit
                  shape = (number of orders x number of fit coefficients)
            ass: numpy array (2D), the fit coefficients array for
                  the widths fit
                  shape = (number of orders x number of fit coefficients)
            tilt: numpy array (1D), the tilt angle of each order

:param image: numpy array (2D), the image
:param orderp: numpy array (2D), the order profile image
:param rnum: int, the order number for this iteration
:param kwargs: additional keywords to pass to the extraction wrapper

        - allowed keywords are:

        range1  (defaults to "IC_EXT_RANGE1")
        range2  (defaults to "IC_EXT_RANGE2")
        gain    (defaults to "GAIN")
        sigdet  (defaults to "SIGDET")

:return cent: numpy array (1D), the extracted pixel values,
             size = image.shape[1] (along the order direction)
:return cpt: int, zero in this case
\end{pythondocstring}
\end{minipage}

%----------------------------------------------------------------------------------------

\begin{minipage}{\textwidth}
\subsection{}

Defined in \spirouEXTOR\path{.}

\begin{pythonbox}
from SpirouDRS import spirou
spirouEXTOR.ExtractWeightOrder(pp, loc, image, orderp, rnum, **kwargs)
spirouEXTOR.spirouEXTOR.extract_weight_order(pp, loc, image, orderp, rnum, **kwargs)
\end{pythonbox}

\begin{pythondocstring}
Extract order with weight but without tilt using
spirouEXTOR.extract_wrapper()

:param p: parameter dictionary, ParamDict containing constants
    Must contain at least:
            IC_EXT_RANGE: float, the upper and lower edge of the order
                          in rows (y-axis) - half-zone width
            gain: float, the gain of the image
            sigdet: float, the read noise of the image

:param loc: parameter dictionary, ParamDict containing data
        Must contain at least:
            acc: numpy array (2D), the fit coefficients array for
                  the centers fit
                  shape = (number of orders x number of fit coefficients)
            ass: numpy array (2D), the fit coefficients array for
                  the widths fit
                  shape = (number of orders x number of fit coefficients)

:param image: numpy array (2D), the image
:param orderp: numpy array (2D), the order profile image
:param rnum: int, the order number for this iteration
:param kwargs: additional keywords to pass to the extraction wrapper

        - allowed keywords are:

        range1  (defaults to "IC_EXT_RANGE")
        range2  (defaults to "IC_EXT_RANGE")
        gain    (defaults to "GAIN")
        sigdet  (defaults to "SIGDET")

:return cent: numpy array (1D), the extracted pixel values,
             size = image.shape[1] (along the order direction)
:return cpt: int, zero in this case
\end{pythondocstring}
\end{minipage}

%----------------------------------------------------------------------------------------




%%%%%%%%%%%%%%%%%%%%%%%%%%%%%%%%%%%%%%%%%%%%%%%%%%%%%%%%
%%
\clearpage
\newpage
\begin{minipage}{\textwidth}
\section{The spirouFLAT module}
\label{ch:the_module:spirouFLAT}
%%
%%%%%%%%%%%%%%%%%%%%%%%%%%%%%%%%%%%%%%%%%%%%%%%%%%%%%%%%

%----------------------------------------------------------------------------------------
\subsection{MeasureBlazeForOrder}

Defined in \spirouFLAT\path{.spirouFLAT.measure_blaze_for_order}

\begin{pythonbox}
from SpirouDRS import spirouFLAT
spirouFLAT.MeasureBlazeForOrder
spirouFLAT.spirouFLAT.measure_blaze_for_order
\end{pythonbox}

\begin{pythondocstring}
Measure the blaze function (for good pixels this is a polynomial fit of
order = fitdegree, for bad pixels = 1.0).

bad pixels are defined as less than or equal to zero

:param y: numpy array (1D), the extracted pixels for this order
:param fitdegree: int, the polynomial degree

:return blaze: numpy array (1D), size = len(y), the blaze function: for
               good pixels this is the value of the fit, for bad pixels the
               value = 1.0
\end{pythondocstring}
\end{minipage}

%----------------------------------------------------------------------------------------



%%%%%%%%%%%%%%%%%%%%%%%%%%%%%%%%%%%%%%%%%%%%%%%%%%%%%%%%
%%
\clearpage
\newpage
\noindent\begin{minipage}{\textwidth}
\section{The spirouImage module}
\label{ch:the_module:spirouImage}
%%
%%%%%%%%%%%%%%%%%%%%%%%%%%%%%%%%%%%%%%%%%%%%%%%%%%%%%%%%

%----------------------------------------------------------------------------------------
\subsection{AddKey}

Defined in \spirouImage\path{.}

\begin{pythonbox}
from SpirouDRS import spirouImage
spirouImage.AddKey(hdict=None, keywordstore=None, value=None)
spirouImage.spirouFITS.add_new_key(hdict=None, keywordstore=None, value=None)
\end{pythonbox}

\begin{pythondocstring}
Add a new key to hdict from keywordstore, if value is not None then the
keywordstore value is updated. Each keywordstore is in form:
        [key, value, comment]    where key and comment are strings
if hdict is None creates a new dictionary

:param hdict: dictionary or None, storage for adding to FITS rec
:param keywordstore: list, keyword list (defined in spirouKeywords.py)
                     must be in form [string, value, string]
:param value: object or None, if any python object (other than None) will
              replace the value in keywordstore (i.e. keywordstore[1]) with
              value, if None uses the value = keywordstore[1]

:return hdict: dictionary, storage for adding to FITS rec
\end{pythondocstring}
\end{minipage}

%----------------------------------------------------------------------------------------
\noindent\begin{minipage}{\textwidth}
\subsection{AddKey1DList}

Defined in \spirouImage\path{.pirouFITS.add_key_1d_list}

\begin{pythonbox}
from SpirouDRS import spirouImage
spirouImage.AddKey1DList(hdict, keywordstore, values=None, dim1name='order')
spirouImage.pirouFITS.add_key_1d_list(hdict, keywordstore, values=None, dim1name='order')
\end{pythonbox}

\begin{pythondocstring}
Add a new 1d list to key using the keywordstorage[0] as prefix in form
keyword = kewordstoreage + row number

:param hdict: dictionary, storage for adding to FITS rec
:param keywordstore: list, keyword list (defined in spirouKeywords.py)
                     must be in form [string, value, string]
:param values: numpy array or 1D list of keys or None

              if numpy array or 1D list will create a set of keys in form
              keyword = kewordstoreage + row number
              where row number is the position in values
              with value = values[row number][column number]

              if None uses the value = keywordstore[1]
:param dim1name: string, the name for dimension 1 (rows), used in FITS rec
                 HEADER comments in form:
      COMMENT = keywordstore[2] dim1name={row number}

:return hdict: dictionary, storage for adding to FITS rec
\end{pythondocstring}
\end{minipage}

%----------------------------------------------------------------------------------------
\noindent\begin{minipage}{\textwidth}
\subsection{AddKey2DList}

Defined in \spirouImage\path{.spirouFITS.add_key_2d_list}

\begin{pythonbox}
from SpirouDRS import spirouImage
spirouImage.AddKey2DList
spirouImage.spirouFITS.add_key_2d_list
\end{pythonbox}

\begin{pythondocstring}
Add a new 2d list to key using the keywordstorage[0] as prefix in form
keyword = kewordstoreage + number

where number = (row number * number of columns) + column number

:param hdict: dictionary, storage for adding to FITS rec
:param keywordstore: list, keyword list (defined in spirouKeywords.py)
                     must be in form [string, value, string]
:param values: numpy array or 2D list of keys or None

              if numpy array or 2D list will create a set of keys in form
              keyword = kewordstoreage + number
              where number = (row number*number of columns)+column number
              with value = values[row number][column number]

              if None uses the value = keywordstore[1]
:param dim1name: string, the name for dimension 1 (rows), used in FITS rec
                 HEADER comments in form:
      COMMENT = keywordstore[2] dim1name={row number} dim2name={col number}
:param dim2name: string, the name for dimension 2 (cols), used in FITS rec
                 HEADER comments in form:
      COMMENT = keywordstore[2] dim1name={row number} dim2name={col number}

:return hdict: dictionary, storage for adding to FITS rec
\end{pythondocstring}
\end{minipage}

%----------------------------------------------------------------------------------------

\noindent\begin{minipage}{\textwidth}
\subsection{ConvertToE}

Defined in \spirouImage\path{.}

\begin{pythonbox}
from SpirouDRS import spirouImage
spirouImage.ConvertToE(image, p=None, gain=None, exptime=None)
spirouImage.spirouImage.convert_to_e(image, p=None, gain=None, exptime=None)
\end{pythonbox}

\begin{pythondocstring}
Converts image from ADU/s into e-

:param image:
:param p: parameter dictionary, ParamDict containing constants
        Must contain at least: (if exptime is None)
            exptime: float, the exposure time of the image
            gain: float, the gain of the image

:param gain: float, if p is None, used as the gain to multiple the image by
:param exptime: float, if p is None, used as the exposure time the image
                is multiplied by

:return newimage: numpy array (2D), the image in e-
\end{pythondocstring}
\end{minipage}

%----------------------------------------------------------------------------------------
\noindent\begin{minipage}{\textwidth}
\subsection{ConvertToADU}

Defined in \spirouImage\path{.spirouImage.convert_to_adu}

\begin{pythonbox}
from SpirouDRS import spirouImage
spirouImage.ConvertToADU(image, p=None, exptime=None)
spirouImage.spirouImage.convert_to_adu(image, p=None, exptime=None)
\end{pythonbox}

\begin{pythondocstring}
Converts image from ADU/s into ADU

:param image:

:param p: parameter dictionary, ParamDict containing constants
    Must contain at least: (if exptime is None)
        exptime: float, the exposure time of the image

:param exptime: float, if p is None, used as the exposure time the image
                is multiplied by

:return newimage: numpy array (2D), the image in e-
\end{pythondocstring}
\end{minipage}

%----------------------------------------------------------------------------------------

\noindent\begin{minipage}{\textwidth}
\subsection{CopyOriginalKeys}

Defined in \spirouImage\path{.spirouFITS.copy_original_keys}

\begin{pythonbox}
from SpirouDRS import spirouImage
spirouImage.CopyOriginalKeys(header, comments, hdict=None, forbid_keys=True)
spirouImage.spirouFITS.copy_original_keys(header, comments, hdict=None, forbid_keys=True)
\end{pythonbox}

\begin{pythondocstring}
Copies keys from hdr dictionary to hdict, if forbid_keys is True some
keys will not be copies (defined in python code)

:param header: header dictionary from readimage (ReadImage) function

:param comments: comment dictionary from readimage (ReadImage) function

:param hdict: dictionary or None, header dictionary to write to fits file
              if None hdict is created

            Must be in form:

                    hdict[key] = (value, comment)
            or
                    hdict[key] = value     (comment will be equal to
                                            "UNKNOWN"
:param forbid_keys: bool, if True uses the forbidden copy keys (defined in
                    spirouConfig.Constants.FORBIDDEN_COPY_KEYS() to remove
                    certain keys from those being copied, if False copies
                    all keys from input header

:return hdict: dictionary, (updated or new) header dictionary containing
               key/value pairs from the header (that are NOT in
               spirouConfig.spirouConst.FORBIDDEN_COPY_KEY)
\end{pythondocstring}
\end{minipage}

%----------------------------------------------------------------------------------------

\noindent\begin{minipage}{\textwidth}
\subsection{CopyRootKeys}

Defined in \spirouImage\path{.spirouFITS.copy_root_keys}

\begin{pythonbox}
from SpirouDRS import spirouImage
spirouImage.CopyRootKeys(hdict=None, filename=None, root=None, ext=0)
spirouImage.spirouFITS.copy_root_keys(hdict=None, filename=None, root=None, ext=0)
\end{pythonbox}

\begin{pythondocstring}
Copy keys from a filename to hdict

:param hdict: dictionary or None, header dictionary to write to fits file
              if None hdict is created
:param filename: string, location and filename of the FITS rec to open

:param ext: int, the extension of the FITS rec to open header from
            (defaults to 0)
:return:
\end{pythondocstring}
\end{minipage}

%----------------------------------------------------------------------------------------

\noindent\begin{minipage}{\textwidth}
\subsection{CorrectForDark}

Defined in \spirouImage\path{.spirouImage.correct_for_dark}

\begin{pythonbox}
from SpirouDRS import spirouImage
spirouImage.CorrectForDark(p, image, header, nfiles=None, return_dark=False)
spirouImage.spirouImage.correct_for_dark(p, image, header, nfiles=None, return_dark=False)
\end{pythonbox}

\begin{pythondocstring}
Corrects "image" for "dark" using calibDB file (header must contain
value of p['ACQTIME_KEY'] as a keyword)

:param p: parameter dictionary, ParamDict containing constants
    Must contain at least:
            nbframes: int, the number of frames/files (usually the length
                      of "arg_file_names")
            calibDB: dictionary, the calibration database dictionary
                     (if not in "p" we construct it and need "max_time_unix"
            max_time_unix: float, the unix time to use as the time of
                            reference (used only if calibDB is not defined)
            log_opt: string, log option, normally the program name
            DRS_CALIB_DB: string, the directory that the calibration
                          files should be saved to/read from

:param image: numpy array (2D), the image
:param header: dictionary, the header dictionary created by
               spirouFITS.ReadImage
:param nfiles: int or None, number of files that created image (need to
               multiply by this to get the total dark) if None uses
               p['nbframes']
:param return_dark: bool, if True returns corrected_image and dark
                    if False (default) returns corrected_image

:return corrected_image: numpy array (2D), the dark corrected image
                         only returned if return_dark = True:
:return darkimage: numpy array (2D), the dark
\end{pythondocstring}
\end{minipage}

%----------------------------------------------------------------------------------------

\noindent\begin{minipage}{\textwidth}
\subsection{FitTilt}

Defined in \spirouImage\path{.spirouImage.fit_tilt}

\begin{pythonbox}
from SpirouDRS import spirouImage
spirouImage.FitTilt(pp, lloc)
spirouImage.spirouImage.fit_tilt(pp, lloc)
\end{pythonbox}

\begin{pythondocstring}
Fit the tilt (lloc['tilt'] with a polynomial of size = p['ic_tilt_filt']
return the coefficients, fit and residual rms in lloc dictionary

:param pp: parameter dictionary, ParamDict containing constants
    Must contain at least:
        IC_TILT_FIT: int, Order of polynomial to fit for tilt

:param loc: parameter dictionary, ParamDict containing data
        Must contain at least:
            number_orders: int, the number of orders in reference spectrum
            tilt: numpy array (1D), the tilt angle of each order

:return loc: parameter dictionary, the updated parameter dictionary
        Adds/updates the following:
            xfit_tilt: numpy array (1D), the order numbers
            yfit_tilt: numpy array (1D), the fit for the tilt angle of each
                       order
            a_tilt: numpy array (1D), the fit coefficients (generated by
                    numpy.polyfit but IN REVERSE ORDER)
            rms_tilt: float, the RMS (np.std) of the residuals of the
                      tilt - tilt fit values
\end{pythondocstring}
\end{minipage}

%----------------------------------------------------------------------------------------

\noindent\begin{minipage}{\textwidth}
\subsection{FlipImage}

Defined in \spirouImage\path{.spirouImage.flip_image}

\begin{pythonbox}
from SpirouDRS import spirouImage
spirouImage.FlipImage(image, fliprows=True, flipcols=True)
spirouImage.spirouImage.flip_image(image, fliprows=True, flipcols=True)
\end{pythonbox}

\begin{pythondocstring}
Flips the image in the x and/or the y direction

:param image: numpy array (2D), the image
:param fliprows: bool, if True reverses row order (axis = 0)
:param flipcols: bool, if True reverses column order (axis = 1)

:return newimage: numpy array (2D), the flipped image
\end{pythondocstring}
\end{minipage}

%----------------------------------------------------------------------------------------

\noindent\begin{minipage}{\textwidth}
\subsection{GetAllSimilarFiles}

Defined in \spirouImage\path{.spirouImage.get_all_similar_files}

\begin{pythonbox}
from SpirouDRS import spirouImage
spirouImage.GetAllSimilarFiles(p, directory, prefix=None, suffix=None)
spirouImage.spirouImage.get_all_similar_files(p, directory, prefix=None, suffix=None)
\end{pythonbox}

\begin{pythondocstring}
Get all similar files in a directory with matching prefix and suffix defined
either by "prefix" and "suffix" or by p["ARG_FILE_NAMES"][0]

:param p: parameter dictionary, ParamDict containing constants
    Must contain at least:
            arg_file_names: list, list of files taken from the command line
                            (or call to recipe function) must have at least
                            one string filename in the list
            log_opt: string, log option, normally the program name

:param directory: string, the directory to search for files
:param prefix: string or None, if not None the prefix to search for, if
               None defines the prefix from the first 5 characters of
               p["ARG_FILE_NAMES"][0]
:param suffix: string  or None, if not None the suffix to search for, if
               None defines the prefix from the last 8 characters of
               p["ARG_FILE_NAMES"][0]

:return filelist: list of strings, the full paths of all files that are in
                  "directory" with the matching prefix and suffix defined
                  either by "prefix" and "suffix" or by
                  p["ARG_FILE_NAMES"][0]
\end{pythondocstring}
\end{minipage}

%----------------------------------------------------------------------------------------

\noindent\begin{minipage}{\textwidth}
\subsection{GetSigdet}

Defined in \spirouImage\path{.spirouImage.get_sigdet}

\begin{pythonbox}
from SpirouDRS import spirouImage
spirouImage.GetSigdet(p, hdr, name=None, return_value=False)
spirouImage.spirouImage.get_sigdet(p, hdr, name=None, return_value=False)
\end{pythonbox}

\begin{pythondocstring}
Get sigdet from HEADER. Wrapper for spirouImage.get_param

:param p: parameter dictionary, ParamDict of constants
:param hdr: dictionary, header dictionary to extract
:param name: string or None, if not None the name for the parameter
             logged if there is an error in getting parameter, if name is
             None the name is taken as "keyword"
:param return_value: bool, if True returns parameter, if False adds
                     parameter to "p" parameter dictionary (and sets source)

:return value: if return_value is True value of parameter is returned
:return p: if return_value is False, updated parameter dictionary p with
           key = name is returned
\end{pythondocstring}
\end{minipage}

%----------------------------------------------------------------------------------------

\noindent\begin{minipage}{\textwidth}
\subsection{GetExpTime}

Defined in \spirouImage\path{.spirouImage.get_exptime}

\begin{pythonbox}
from SpirouDRS import spirouImage
spirouImage.GetExpTime(p, hdr, name=None, return_value=False)
spirouImage.spirouImage.get_exptime(p, hdr, name=None, return_value=False)
\end{pythonbox}

\begin{pythondocstring}
Get Exposure time from HEADER. Wrapper for spirouImage.get_param

:param p: parameter dictionary, ParamDict of constants
:param hdr: dictionary, header dictionary to extract
:param name: string or None, if not None the name for the parameter
             logged if there is an error in getting parameter, if name is
             None the name is taken as "keyword"
:param return_value: bool, if True returns parameter, if False adds
                     parameter to "p" parameter dictionary (and sets source)

:return value: if return_value is True value of parameter is returned
:return p: if return_value is False, updated parameter dictionary p with
           key = name is returned
\end{pythondocstring}
\end{minipage}

%----------------------------------------------------------------------------------------

\noindent\begin{minipage}{\textwidth}
\subsection{GetGain}

Defined in \spirouImage\path{.spirouImage.get_gain}

\begin{pythonbox}
from SpirouDRS import spirouImage
spirouImage.GetGain(p, hdr, name=None, return_value=False)
spirouImage.spirouImage.get_gain(p, hdr, name=None, return_value=False)
\end{pythonbox}

\begin{pythondocstring}
Get Gain from HEADER. Wrapper for spirouImage.get_param

:param p: parameter dictionary, ParamDict of constants
:param hdr: dictionary, header dictionary to extract
:param name: string or None, if not None the name for the parameter
             logged if there is an error in getting parameter, if name is
             None the name is taken as "keyword"
:param return_value: bool, if True returns parameter, if False adds
                     parameter to "p" parameter dictionary (and sets source)

:return value: if return_value is True value of parameter is returned
:return p: if return_value is False, updated parameter dictionary p with
           key = name is returned
\end{pythondocstring}
\end{minipage}

%----------------------------------------------------------------------------------------

\noindent\begin{minipage}{\textwidth}
\subsection{GetAcqTime}

Defined in \spirouImage\path{.spirouImage.get_acqtime}

\begin{pythonbox}
from SpirouDRS import spirouImage
spirouImage.GetAcqTime(p, hdr, name=None, kind='human', return_value=False)
spirouImage.spirouImage.get_acqtime(p, hdr, name=None, kind='human', return_value=False)
\end{pythonbox}

\begin{pythondocstring}
Get the acquision time from the header file, if there is not header file
use the parameter dictionary "p" to open the header in 'arg_file_names[0]'

:param p: parameter dictionary, ParamDict containing constants
    Must contain at least:
        "name" defined in call
        parameter dictionary to give to value

:param hdr: dictionary, the header dictionary created by
            spirouFITS.ReadImage
:param name: string, the name in parameter dictionary to give to value
             if return_value is False (i.e. p[name] = value)
:param kind: string, 'human' for 'YYYY-mm-dd-HH-MM-SS.ss' or 'unix'
             for time since 1970-01-01
:param return_value: bool, if False value is returned in p as p[name]
                     if True value is returned

:return p or value: dictionary or string or float, if return_value is False
                    parameter dictionary is returned, if return_value is
                    True and kind=='human' returns a string, if return_value
                    is True and kind=='unix' returns a float
\end{pythondocstring}
\end{minipage}

%----------------------------------------------------------------------------------------

\noindent\begin{minipage}{\textwidth}
\subsection{ReadParam}

Defined in \spirouImage\path{.spirouImage.get_param}

\begin{pythonbox}
from SpirouDRS import spirouImage
spirouImage.ReadParam(p, hdr, name=None, kind='human', return_value=False)
spirouImage.spirouImage.get_param(p, hdr, name=None, kind='human', return_value=False)
\end{pythonbox}

\begin{pythondocstring}
Get the acquision time from the header file, if there is not header file
use the parameter dictionary "p" to open the header in 'arg_file_names[0]'

:param p: parameter dictionary, ParamDict containing constants
    Must contain at least:
        "name" defined in call
        parameter dictionary to give to value

:param hdr: dictionary, the header dictionary created by
            spirouFITS.ReadImage
:param name: string, the name in parameter dictionary to give to value
             if return_value is False (i.e. p[name] = value)
:param kind: string, 'human' for 'YYYY-mm-dd-HH-MM-SS.ss' or 'unix'
             for time since 1970-01-01
:param return_value: bool, if False value is returned in p as p[name]
                     if True value is returned

:return p or value: dictionary or string or float, if return_value is False
                    parameter dictionary is returned, if return_value is
                    True and kind=='human' returns a string, if return_value
                    is True and kind=='unix' returns a float
\end{pythondocstring}
\end{minipage}

%----------------------------------------------------------------------------------------

\noindent\begin{minipage}{\textwidth}
\subsection{GetKey}

Defined in \spirouImage\path{.spirouFITS.keylookup}

\begin{pythonbox}
from SpirouDRS import spirouImage
spirouImage.GetKey(p, d=None, key=None, has_default=False, default=None)
spirouImage.spirouFITS.keylookup(p, d=None, key=None, has_default=False, default=None)
\end{pythonbox}

\begin{pythondocstring}
Looks for a key in dictionary "p" or "d", if has_default is True sets
value of key to 'default' if not found else logs an error

:param p: parameter dictionary, ParamDict containing constants
    Must contain at least:
            log_opt: string, log option, normally the program name
            "key": if d is None must contain key="key" or error is raised
:param d: dictionary, any dictionary, if None uses parameter dictionary
          if "d" is not None then must contain key="key" or error is raised
:param key: string, key in the dictionary to find
:param has_default: bool, if True uses "default" as the value if key
                    not found
:param default: object, value of the key if not found and
                has_default is True

:return value: object, value of p[key] or default (if has_default=True)
\end{pythondocstring}
\end{minipage}

%----------------------------------------------------------------------------------------

\noindent\begin{minipage}{\textwidth}
\subsection{GetKeys}

Defined in \spirouImage\path{.spirouFITS.keyslookup}

\begin{pythonbox}
from SpirouDRS import spirouImage
spirouImage.GetKeys(p, d=None, keys=None, has_default=False, defaults=None)
spirouImage.spirouFITS.keyslookup(p, d=None, keys=None, has_default=False, defaults=None)
\end{pythonbox}

\begin{pythondocstring}
Looks for keys in dictionary "p" or "d", if has_default is True sets
value of key to 'default' if not found else logs an error

:param p: parameter dictionary, ParamDict containing constants
    Must contain at least:
            log_opt: string, log option, normally the program name
            "key": if d is None must contain key="key" or error is raised
:param d: dictionary, any dictionary, if None uses parameter dictionary
          if "d" is not None then must contain key="key" or error is raised
:param keys: list of strings, keys in the dictionary to find
:param has_default: bool, if True uses "default" as the value if key
                    not found
:param defaults: list of objects or None, values of the keys if not
                 found and has_default is True

:return values: list of objects, values of p[key] for key in keys
                or default value for each key (if has_default=True)
\end{pythondocstring}
\end{minipage}

%----------------------------------------------------------------------------------------

\noindent\begin{minipage}{\textwidth}
\subsection{GetTilt}

Defined in \spirouImage\path{.}

\begin{pythonbox}
from SpirouDRS import spirouImage
spirouImage.GetTilt(pp, lloc, image)
spirouImage.spirouImage.get_tiltspirouImage.get_tilt(pp, lloc, image)
\end{pythonbox}

\begin{pythondocstring}
Get the tilt by correlating the extracted fibers

:param pp: parameter dictionary, ParamDict containing constants
    Must contain at least:
            ic_tilt_coi: int, oversampling factor
            log_opt: string, log option, normally the program name

:param loc: parameter dictionary, ParamDict containing data
        Must contain at least:
            number_orders: int, the number of orders in reference spectrum
            cent1: numpy array (2D), the extraction for A, updated is
                   the order "rnum"
            cent2: numpy array (2D), the extraction for B, updated is
                   the order "rnum"
            offset: numpy array (1D), the center values with the
                    offset in 'IC_CENT_COL' added

:param image: numpy array (2D), the image

:return loc: parameter dictionary, the updated parameter dictionary
        Adds/updates the following:
            nbcos: numpy array, zero array  (length of "number_orderes")
            tilt: numpy array (1D), the tilt angle of each order
\end{pythondocstring}
\end{minipage}

%----------------------------------------------------------------------------------------

\noindent\begin{minipage}{\textwidth}
\subsection{GetTypeFromHeader}

Defined in \spirouImage\path{.spirouFITS.get_type_from_header}

\begin{pythonbox}
from SpirouDRS import spirouImage
spirouImage.GetTypeFromHeader(p, keywordstore, hdict=None, filename=None)
spirouImage.spirouFITS.get_type_from_header(p, keywordstore, hdict=None, filename=None)
\end{pythonbox}

\begin{pythondocstring}
Special FITS HEADER keyword - get the type of file from a FITS file HEADER
using "keywordstore"

:param p: parameter dictionary, ParamDict containing constants
    Must contain at least:
            log_opt: string, log option, normally the program name
            fitsfilename: string, the full path of for the main raw fits
                          file for a recipe
                          i.e. /data/raw/20170710/filename.fits

:param keywordstore: list, a keyword store in the form
                     [name, value, comment] where the format is
                     [string, object, string]
:param hdict: dictionary or None, the HEADER dictionary containing
              key/value pairs from a FITS HEADER, if None uses the
              header from "FITSFILENAME" in "p", unless filename is not None
              This hdict is used to get the type of file
:param filename: string or None, if not None and hdict is None, this is the
                 file which is used to extract the HEADER from to get
                 the type of file

:return ftype: string, the type of file (extracted from a HEADER dictionary/
               file) if undefined set to 'UNKNOWN'
\end{pythondocstring}
\end{minipage}

%----------------------------------------------------------------------------------------

\noindent\begin{minipage}{\textwidth}
\subsection{LocateBadPixels}

Defined in \spirouImage\path{.spirouImage.locate_bad_pixels}

\begin{pythonbox}
from SpirouDRS import spirouImage
spirouImage.LocateBadPixels(p, fimage, fmed, dimage, wmed=None)
spirouImage.spirouImage.locate_bad_pixels(p, fimage, fmed, dimage, wmed=None)
\end{pythonbox}

\begin{pythondocstring}
Locate the bad pixels in the flat image and the dark image

:param p: parameter dictionary, ParamDict containing constants
    Must contain at least:
            log_opt: string, log option, normally the program name
            BADPIX_FLAT_MED_WID: float, the median image in the x
                                 dimension over a boxcar of this width
            BADPIX_FLAT_CUT_RATIO: float, the maximum differential pixel
                                   cut ratio
            BADPIX_ILLUM_CUT: float, the illumination cut parameter
            BADPIX_MAX_HOTPIX: float, the maximum flux in ADU/s to be
                               considered too hot to be used

:param fimage: numpy array (2D), the flat normalised image
:param fmed: numpy array (2D), the flat median normalised image
:param dimage: numpy array (2D), the dark image
:param wmed: float or None, if not None defines the median filter width
             if None uses p["BADPIX_MED_WID", see
             scipy.ndimage.filters.median_filter "size" for more details

:return bad_pix_mask: numpy array (2D), the bad pixel mask image
:return badpix_stats: list of floats, the statistics array:
                        Fraction of hot pixels from dark [%]
                        Fraction of bad pixels from flat [%]
                        Fraction of NaN pixels in dark [%]
                        Fraction of NaN pixels in flat [%]
                        Fraction of bad pixels with all criteria [%]
\end{pythondocstring}
\end{minipage}

%----------------------------------------------------------------------------------------

\noindent\begin{minipage}{\textwidth}
\subsection{MakeTable}

Defined in \spirouImage\path{.spirouTable.make_table}

\begin{pythonbox}
from SpirouDRS import spirouImage
spirouImage.MakeTable
spirouImage.spirouTable.make_table
\end{pythonbox}

\begin{pythondocstring}
Construct an astropy table from columns and values

:param columns: list of strings, the list of column names
:param values: list of lists or numpy array (2D), the list of lists/array
               of values, first dimension must have same length as number
               of columns, there must be the same number of values in each
               column
:param formats: list of strings, the astropy formats for each column
                i.e. 0.2f  for a float with two decimal places, must have
                same length as number of columns
:param units: list of strings, the units for each column, must have
              same length as number of columns

:return table: astropy.table.Table instance, the astropy table containing
               all columns and data
\end{pythondocstring}
\end{minipage}

%----------------------------------------------------------------------------------------

\noindent\begin{minipage}{\textwidth}
\subsection{MeasureDark}

Defined in \spirouImage\path{.spirouImage.measure_dark}

\begin{pythonbox}
from SpirouDRS import spirouImage
spirouImage.MeasureDark(pp, image, image_name, short_name)
spirouImage.spirouImage.measure_dark(pp, image, image_name, short_name)
\end{pythonbox}

\begin{pythondocstring}
Measure the dark pixels in "image"

:param pp: parameter dictionary, ParamDict containing constants
    Must contain at least:
            log_opt: string, log option, normally the program name
            DARK_QMIN: int, The lower percentile (0 - 100)
            DARK_QMAX: int, The upper percentile (0 - 100)
            HISTO_BINS: int,  The number of bins in dark histogram
            HISTO_RANGE_LOW: float, the lower extent of the histogram
                             in ADU/s
            HISTO_RANGE_HIGH: float, the upper extent of the histogram
                              in ADU/s

:param image: numpy array (2D), the image
:param image_name: string, the name of the image (for logging)
:param short_name: string, suffix (for parameter naming -
                    parmaeters added to pp with suffix i)

:return pp: parameter dictionary, the updated parameter dictionary
        Adds the following: (based on "short_name")
            histo_full: numpy.histogram tuple (hist, bin_edges) for
                        the full image
            histo_blue: numpy.histogram tuple (hist, bin_edges) for
                        the blue part of the image
            histo_red: numpy.histogram tuple (hist, bin_edges) for
                        the red part of the image
            med_full: float, the median value of the non-Nan image values
                      for the full image
            med_blue: float, the median value of the non-Nan image values
                      for the blue part of the image
            med_red: float, the median value of the non-Nan image values
                     for the red part of the image
            dadead_full: float, the fraction of dead pixels as a percentage
                         for the full image
            dadead_blue: float, the fraction of dead pixels as a percentage
                         for the blue part of the image
            dadead_red: float, the fraction of dead pixels as a percentage
                        for the red part of the image

      where:
          hist : numpy array (1D) The values of the histogram.
          bin_edges : numpy array (1D) of floats, the bin edges
\end{pythondocstring}
\end{minipage}

%----------------------------------------------------------------------------------------

\noindent\begin{minipage}{\textwidth}
\subsection{NormMedianFlat}

Defined in \spirouImage\path{.spirouImage.normalise_median_flat}

\begin{pythonbox}
from SpirouDRS import spirouImage
spirouImage.NormMedianFlat(p, image, method='new', wmed=None, percentile=None)
spirouImage.spirouImage.normalise_median_flat(p, image, method='new', wmed=None, percentile=None)
\end{pythonbox}

\begin{pythondocstring}
Applies a median filter and normalises. Median filter is applied with width
"wmed" or p["BADPIX_FLAT_MED_WID"] if wmed is None) and then normalising by
the 90th percentile

:param p: parameter dictionary, ParamDict containing constants
    Must contain at least:
            BADPIX_FLAT_MED_WID: float, the median image in the x
                                 dimension over a boxcar of this width
            BADPIX_NORM_PERCENTILE: float, the percentile to normalise
                                    to when normalising and median
                                    filtering image
            log_opt: string, log option, normally the program name

:param image: numpy array (2D), the iamge to median filter and normalise
:param method: string, "new" or "old" if "new" uses np.percentile else
               sorts the flattened image and takes the "percentile" (i.e.
               90th) pixel value to normalise
:param wmed: float or None, if not None defines the median filter width
             if None uses p["BADPIX_MED_WID", see
             scipy.ndimage.filters.median_filter "size" for more details
:param percentile: float or None, if not None degines the percentile to
                   normalise the image at, if None used from
                   p["BADPIX_NORM_PERCENTILE"]

:return norm_med_image: numpy array (2D), the median filtered and normalised
                        image
:return norm_image: numpy array (2D), the normalised image
\end{pythondocstring}
\end{minipage}

%----------------------------------------------------------------------------------------

\noindent\begin{minipage}{\textwidth}
\subsection{ReadData}

Defined in \spirouImage\path{.spirouFITS.readdata}

\begin{pythonbox}
from SpirouDRS import spirouImage
spirouImage.ReadData(p, filename, log=True)
spirouImage.spirouFITS.readdata(p, filename, log=True)
\end{pythonbox}

\begin{pythondocstring}
Reads the image 'fitsfilename' defined in p and adds files defined in
'arg_file_names' if add is True

:param p: parameter dictionary, ParamDict containing constants
    Must contain at least:
            log_opt: string, log option, normally the program name

:param filename: string, filename of the image to read
:param log: bool, if True logs opening and size

:return image: numpy array (2D), the image
:return header: dictionary, the header file of the image
:return nx: int, the shape in the first dimension, i.e. data.shape[0]
:return ny: int, the shape in the second dimension, i.e. data.shape[1]
\end{pythondocstring}
\end{minipage}

%----------------------------------------------------------------------------------------

\noindent\begin{minipage}{\textwidth}
\subsection{ReadImage}

Defined in \spirouImage\path{.spirouFITS.readimage}

\begin{pythonbox}
from SpirouDRS import spirouImage
spirouImage.ReadImage(p, filename=None, log=True, kind=None)
spirouImage.spirouFITS.readimage(p, filename=None, log=True, kind=None)
\end{pythonbox}

\begin{pythondocstring}
Reads the image 'fitsfilename' defined in p and adds files defined in
'arg_file_names' if add is True

:param p: parameter dictionary, ParamDict containing constants
    Must contain at least:
            fitsfilename: string, the full path of for the main raw fits
                  file for a recipe
                  i.e. /data/raw/20170710/filename.fits
            log_opt: string, log option, normally the program name
            arg_file_names: list, list of files taken from the command line
                            (or call to recipe function) must have at least
                            one string filename in the list

:param filename: string or None, filename of the image to read, if None
                 then p['fitsfilename'] is used
:param log: bool, if True logs opening and size
:param kind: string or None, if defined names the image else just image,
             used in logging (if log = True)

:return image: numpy array (2D), the image
:return header: dictionary, the header file of the image
:return nx: int, the shape in the first dimension, i.e. data.shape[0]
:return ny: int, the shape in the second dimension, i.e. data.shape[1]
\end{pythondocstring}
\end{minipage}

%----------------------------------------------------------------------------------------

\noindent\begin{minipage}{\textwidth}
\subsection{ReadTable}

Defined in \spirouImage\path{.spirouTable.read_table}

\begin{pythonbox}
from SpirouDRS import spirouImage
spirouImage.ReadTable(filename, fmt, colnames=None)
spirouImage.spirouTable.read_table(filename, fmt, colnames=None)
\end{pythonbox}

\begin{pythondocstring}
Reads a table from file "filename" in format "fmt", if colnames are defined
renames the columns to these name

:param filename: string, the filename and location of the table to read
:param fmt: string, the format of the table to read from (must be valid
            for astropy.table to read - see below)
:param colnames: list of strings or None, if not None renames all columns
                 to these strings, must be the same length as columns
                 in file that is read

:return None:

astropy.table readable formats are as follows:
\end{pythondocstring}
\end{minipage}

%----------------------------------------------------------------------------------------

\noindent\begin{minipage}{\textwidth}
\subsection{ReadImageAndCombine}

Defined in \spirouImage\path{.spirouFITS.readimage_and_combine}

\begin{pythonbox}
from SpirouDRS import spirouImage
spirouImage.ReadImageAndCombine(p, framemath='+', filename=None, filenames=None, log=True)
spirouImage.spirouFITS.readimage_and_combine(p, framemath='+', filename=None, filenames=None, log=True)
\end{pythonbox}

\begin{pythondocstring}
Reads the image 'fitsfilename' defined in p and adds files defined in
'arg_file_names' if add is True

:param p: parameter dictionary, ParamDict containing constants
    Must contain at least:
            log_opt: string, log option, normally the program name
            
            optional:
            fitsfilename: string, the full path of for the main raw fits
                  file for a recipe i.e. /data/raw/20170710/filename.fits
                  (if filename is None this is required)
                  
            arg_file_names: list, list of files taken from the command line
                            (or call to recipe function) must have at least
                            one string filename in the list
                  (if filenames is None this is required)

:param framemath: string, controls how files should be added

            currently supported are:
            'add' or '+'           - adds the frames
            'sub' or '-'           - subtracts the frames
            'average' or 'mean'    - averages the frames
            'multiply' or '*'      - multiplies the frames
            'divide' or '/'        - divides the frames
            'none'                 - does not add
            
:param filename: string or None, filename of the image to read, if None
                 then p['fitsfilename'] is used

:param filenames: list of strings or None, filenames to combine with 
                  "filename", if None then p['arg_file_names'] is used

:param log: bool, if True logs opening and size

:return image: numpy array (2D), the image
:return header: dictionary, the header file of the image
:return nx: int, the shape in the first dimension, i.e. data.shape[0]
:return ny: int, the shape in the second dimension, i.e. data.shape[1]
\end{pythondocstring}
\end{minipage}

%----------------------------------------------------------------------------------------

\noindent\begin{minipage}{\textwidth}
\subsection{ReadFlatFile}

Defined in \spirouImage\path{.spirouFITS.read_flat_file}

\begin{pythonbox}
from SpirouDRS import spirouImage
spirouImage.ReadFlatFile(p, hdr=None, filename=None, key=None)
spirouImage.spirouFITS.read_flat_file(p, hdr=None, filename=None, key=None)
\end{pythonbox}

\begin{pythondocstring}
Reads the wave file (from calib database or filename)

:param p: parameter dictionary, ParamDict containing constants
    Must contain at least:
            fitsfilename: string, the full path of for the main raw fits
                          file for a recipe
                          i.e. /data/raw/20170710/filename.fits
            fiber: string, the fiber used for this recipe (eg. AB or A or C)
            log_opt: string, log option, normally the program name

:param hdr: dictionary or None, the header dictionary to look for the
                 acquisition time in, if None loads the header from
                 p['fitsfilename']
:param filename: string or None, the filename and path of the tilt file,
                 if None gets the TILT file from the calib database
                 keyword "TILT"
:param key: string or None, if None key='WAVE' else uses string as key
            from calibDB (first entry) to get wave file

:return wave: list of the tilt for each order
\end{pythondocstring}
\end{minipage}

%----------------------------------------------------------------------------------------

\noindent\begin{minipage}{\textwidth}
\subsection{ReadHeader}

Defined in \spirouImage\path{.spirouFITS.read_header}

\begin{pythonbox}
from SpirouDRS import spirouImage
spirouImage.ReadHeader(p=None, filepath=None, ext=0)
spirouImage.spirouFITS.read_header(p=None, filepath=None, ext=0)
\end{pythonbox}

\begin{pythondocstring}
Read the header from a file at "filepath" with extention "ext" (default=0)

:param p: parameter dictionary, ParamDict containing constants
    Must contain at least:
            log_opt: string, log option, normally the program name

:param filepath: string, filename and path of FITS file to open
:param ext: int, extension in FITS rec to open (default = 0)

:return hdict: dictionary, the dictionary with key value pairs
\end{pythondocstring}
\end{minipage}

%----------------------------------------------------------------------------------------

\noindent\begin{minipage}{\textwidth}
\subsection{ReadKey}

Defined in \spirouImage\path{.spirouFITS.read_key}

\begin{pythonbox}
from SpirouDRS import spirouImage
spirouImage.ReadKey(p, hdict=None, key=None)
spirouImage.spirouFITS.read_key(p, hdict=None, key=None)
\end{pythonbox}

\begin{pythondocstring}
Read a key from hdict (or p if hdict is not defined) and return it's
value.

:param p: parameter dictionary, ParamDict containing constants
    Must contain at least:
            log_opt: string, log option, normally the program name

:param hdict: dictionary or None, the dictionary to add the key to once
              found, if None creates a new dictionary
:param key: string, key in the dictionary to find

:return value: object, the value of the key from hdict
               (or p if hdict is None)
\end{pythondocstring}
\end{minipage}

%----------------------------------------------------------------------------------------

\noindent\begin{minipage}{\textwidth}
\subsection{Read2Dkey}

Defined in \spirouImage\path{.spirouFITS.read_key_2d_list}

\begin{pythonbox}
from SpirouDRS import spirouImage
spirouImage.Read2Dkey(p, hdict, key, dim1, dim2)
spirouImage.spirouFITS.read_key_2d_list(p, hdict, key, dim1, dim2)
\end{pythonbox}

\begin{pythondocstring}
Read a set of header keys that were created from a 2D list

:param p: parameter dictionary, ParamDict containing constants
    Must contain at least:
            log_opt: string, log option, normally the program name

:param hdict: dictionary, HEADER dictionary to extract key/value pairs from
:param key: string, prefix of HEADER key to construct 2D list from
             key[number]

             where number = (row number * number of columns) + column number
             where column number = dim2 and row number = range(0, dim1)
:param dim1: int, the number of elements in dimension 1 (number of rows)
:param dim2: int, the number of columns in dimension 2 (number of columns)

:return value: numpy array (2D), the reconstructed 2D list of variables
               from the HEADER dictionary keys
\end{pythondocstring}
\end{minipage}

%----------------------------------------------------------------------------------------

\noindent\begin{minipage}{\textwidth}
\subsection{ReadTiltFile}

Defined in \spirouImage\path{.spirouFITS.read_tilt_file}

\begin{pythonbox}
from SpirouDRS import spirouImage
spirouImage.ReadTiltFile(p, hdr=None, filename=None, key=None)
spirouImage.spirouFITS.read_tilt_file(p, hdr=None, filename=None, key=None)
\end{pythonbox}

\begin{pythondocstring}
Reads the tilt file (from calib database or filename) and using the
'kw_TILT' keyword-store extracts the tilts for each order

:param p: parameter dictionary, ParamDict containing constants
    Must contain at least:
            fitsfilename: string, the full path of for the main raw fits
                          file for a recipe
                          i.e. /data/raw/20170710/filename.fits
            kw_TILT: list, the keyword list for kw_TILT (defined in
                     spirouKeywords.py)
            IC_TILT_NBO: int, Number of orders in tilt file

:param hdr: dictionary or None, the header dictionary to look for the
                 acquisition time in, if None loads the header from
                 p['fitsfilename']
:param filename: string or None, the filename and path of the tilt file,
                 if None gets the TILT file from the calib database
                 keyword "TILT"
:param key: string or None, if None key='TILT' else uses string as key
            from calibDB (first entry) to get tilt file

:return tilt: list of the tilt for each order
\end{pythondocstring}
\end{minipage}

%----------------------------------------------------------------------------------------

\noindent\begin{minipage}{\textwidth}
\subsection{ReadWaveFile}

Defined in \spirouImage\path{.spirouFITS.read_wave_file}

\begin{pythonbox}
from SpirouDRS import spirouImage
spirouImage.ReadWaveFile(p, hdr=None, filename=None, key=None, return_header=False)
spirouImage.spirouFITS.read_wave_file(p, hdr=None, filename=None, key=None, return_header=False)
\end{pythonbox}

\begin{pythondocstring}
Reads the wave file (from calib database or filename)

:param p: parameter dictionary, ParamDict containing constants
    Must contain at least:
            fitsfilename: string, the full path of for the main raw fits
                          file for a recipe
                          i.e. /data/raw/20170710/filename.fits
            fiber: string, the fiber used for this recipe (eg. AB or A or C)

:param hdr: dictionary or None, the header dictionary to look for the
                 acquisition time in, if None loads the header from
                 p['fitsfilename']
:param filename: string or None, the filename and path of the tilt file,
                 if None gets the TILT file from the calib database
                 keyword "TILT"
:param key: string or None, if None key='WAVE' else uses string as key
            from calibDB (first entry) to get wave file

:param return_header: bool, if True returns header file else just returns
                      wave file
:return wave: list of the tilt for each order
\end{pythondocstring}
\end{minipage}

%----------------------------------------------------------------------------------------

\noindent\begin{minipage}{\textwidth}
\subsection{ReadOrderProfile}

Defined in \spirouImage\path{.spirouFITS.read_order_profile_superposition}

\begin{pythonbox}
from SpirouDRS import spirouImage
spirouImage.ReadOrderProfile(p, hdr=None, filename=None)
spirouImage.spirouFITS.read_order_profile_superposition(p, hdr=None, filename=None)
\end{pythonbox}

\begin{pythondocstring}
Read the order profile superposition image from either "filename" (if not
None) or get filename from the calibration database using "p"

"ORDER_PROFILE_{X}" must be in calibration database if filename is None
where X is either p["ORDERP_FILE"] or p["FIBER"] (presedence in that order)

:param p: parameter dictionary, ParamDict containing constants
    Must contain at least:
            ORDERP_FILE: string, the suffix for the order profile
                         calibration database key (usually the fiber type)
                         - read from "orderp_file_fpall"
            fiber: string, the fiber used for this recipe (eg. AB or A or C)
            log_opt: string, log option, normally the program name

:param hdr: dictionary or None, header dictionary (used to get the
            acquisition time if trying to get "ORDER_PROFILE_{X}" from
            the calibration database, if None uses the header from the
            first file in "ARG_FILE_NAMES" i.e. "FITSFILENAME"
:param filename: string or None, if defined no need for "hdr" or keys from
                 "p" the order profile is read straight from "filename"

:return orderp: numpy array (2D), the order profile image read from file
\end{pythondocstring}
\end{minipage}

%----------------------------------------------------------------------------------------

\noindent\begin{minipage}{\textwidth}
\subsection{ResizeImage}

Defined in \spirouImage\path{.spirouImage.resize}

\begin{pythonbox}
from SpirouDRS import spirouImage
spirouImage.ResizeImage(image, x=None, y=None, xlow=0, xhigh=None, ylow=0, yhigh=None, getshape=True)
spirouImage.spirouImage.resize(image, x=None, y=None, xlow=0, xhigh=None, ylow=0, yhigh=None, getshape=True)
\end{pythonbox}

\begin{pythondocstring}
Resize an image based on a pixel values

:param image: numpy array (2D), the image
:param x: None or numpy array (1D), the list of x pixels
:param y: None or numpy array (1D), the list of y pixels
:param xlow: int, x pixel value (x, y) in the bottom left corner,
             default = 0
:param xhigh:  int, x pixel value (x, y) in the top right corner,
             if None default is image.shape(1)
:param ylow: int, y pixel value (x, y) in the bottom left corner,
             default = 0
:param yhigh: int, y pixel value (x, y) in the top right corner,
             if None default is image.shape(0)
:param getshape: bool, if True returns shape of newimage with newimage

if getshape = True
:return newimage: numpy array (2D), the new resized image
:return nx: int, the shape in the first dimension, i.e. data.shape[0]
:return ny: int, the shape in the second dimension, i.e. data.shape[1]

if getshape = False
:return newimage: numpy array (2D), the new resized image
\end{pythondocstring}
\end{minipage}

%----------------------------------------------------------------------------------------

\noindent\begin{minipage}{\textwidth}
\subsection{WriteImage}

Defined in \spirouImage\path{.spirouFITS.writeimage}

\begin{pythonbox}
from SpirouDRS import spirouImage
spirouImage.WriteImage(filename, image, hdict, dtype=None)
spirouImage.spirouFITS.writeimage(filename, image, hdict, dtype=None)
\end{pythonbox}

\begin{pythondocstring}
Writes an image and its header to file

:param filename: string, filename to save the fits file to
:param image: numpy array (2D), the image
:param hdict: dictionary, header dictionary to write to fits file

            Must be in form:

                    hdict[key] = (value, comment)
            or
                    hdict[key] = value     (comment will be equal to
                                            "UNKNOWN"
:param dtype: None or hdu format type, forces the image to be in the
              format type specified (if not None)

              valid formats are for example: 'int32', 'float64'

:return None:
\end{pythondocstring}
\end{minipage}

%----------------------------------------------------------------------------------------

\noindent\begin{minipage}{\textwidth}
\subsection{WriteTable}

Defined in \spirouImage\path{.spirouTable.write_table}

\begin{pythonbox}
from SpirouDRS import spirouImage
spirouImage.WriteTable(table, filename, fmt='fits')
spirouImage.spirouTable.write_table(table, filename, fmt='fits')
\end{pythonbox}

\begin{pythondocstring}
Writes a table to file "filename" with format "fmt"

:param filename: string, the filename and location of the table to read
:param fmt: string, the format of the table to read from (must be valid
            for astropy.table to read - see below)

:return None:
\end{pythondocstring}
\end{minipage}

%----------------------------------------------------------------------------------------




%%%%%%%%%%%%%%%%%%%%%%%%%%%%%%%%%%%%%%%%%%%%%%%%%%%%%%%%
%%
\clearpage
\newpage
\noindent\begin{minipage}{\textwidth}
\section{The spirouLOCOR module}
\label{ch:the_module:spirouLOCOR}
%%
%%%%%%%%%%%%%%%%%%%%%%%%%%%%%%%%%%%%%%%%%%%%%%%%%%%%%%%%

%----------------------------------------------------------------------------------------
\subsection{BoxSmoothedImage}

Defined in \spirouLOCOR\path{.spirouLOCOR.smoothed_boxmean_image}

\begin{pythonbox}
from SpirouDRS import spirouLOCOR
spirouLOCOR.BoxSmoothedImage(image, size, weighted=True, mode='convolve')
spirouLOCOR.spirouLOCOR.smoothed_boxmean_image(image, size, weighted=True, mode='convolve')
\end{pythonbox}

\begin{pythondocstring}
Produce a (box) smoothed image, smoothed by the mean of a box of
    size=2*"size" pixels.

    if mode='convolve' (default) then this is done
    by convolving a top-hat function with the image (FAST)
    - note produces small inconsistencies due to FT of top-hat function

    if mode='manual' then this is done by working out the mean in each
    box manually (SLOW)

:param image: numpy array (2D), the image
:param size: int, the number of pixels to mask before and after pixel
             (for every row)
             i.e. box runs from  "pixel-size" to "pixel+size" unless
             near an edge
:param weighted: bool, if True pixel values less than zero are weighted to
                 a value of 1e-6 and values above 0 are weighted to a value
                 of 1
:param mode: string, if 'convolve' convolves with a top-hat function of the
                     size "box" for each column (FAST) - note produces small
                     inconsistencies due to FT of top-hat function

                     if 'manual' calculates every box individually (SLOW)

:return newimage: numpy array (2D), the smoothed image
\end{pythondocstring}
\end{minipage}

%----------------------------------------------------------------------------------------
\noindent\begin{minipage}{\textwidth}
\subsection{CalcLocoFits}

Defined in \spirouLOCOR\path{.spirouLOCOR.calculate_location_fits}

\begin{pythonbox}
from SpirouDRS import spirouLOCOR
spirouLOCOR.CalcLocoFits(coeffs, dim)
spirouLOCOR.spirouLOCOR.calculate_location_fits(coeffs, dim)
\end{pythonbox}

\begin{pythondocstring}
Calculates all fits in coeffs array across pixels of size=dim

:param coeffs: coefficient array,
               size = (number of orders x number of coefficients in fit)
               output array will be size = (number of orders x dim)
:param dim: int, number of pixels to calculate fit for
            fit will be done over x = 0 to dim in steps of 1
:return yfits: array,
               size = (number of orders x dim)
               the fit for each order at each pixel values from 0 to dim
\end{pythondocstring}
\end{minipage}

%----------------------------------------------------------------------------------------
\noindent\begin{minipage}{\textwidth}
\subsection{FindPosCentCol}

Defined in \spirouLOCOR\path{.spirouLOCOR.find_position_of_cent_col}

\begin{pythonbox}
from SpirouDRS import spirouLOCOR
spirouLOCOR.FindPosCentCol(values, threshold)
spirouLOCOR.spirouLOCOR.find_position_of_cent_col(values, threshold)
\end{pythonbox}

\begin{pythondocstring}
Finds the central positions based on the central column values

:param values: numpy array (1D) size = number of rows,
                the central column values
:param threshold: float, the threshold above which to find pixels as being
                  part of an order

:return position: numpy array (1D), size= number of rows,
                  the pixel positions in cvalues where the centres of each
                  order should be
\end{pythondocstring}
\end{minipage}

%----------------------------------------------------------------------------------------
\noindent\begin{minipage}{\textwidth}
\subsection{FindOrderCtrs}

Defined in \spirouLOCOR\path{.spirouLOCOR.find_order_centers}

\begin{pythonbox}
from SpirouDRS import spirouLOCOR
spirouLOCOR.FindOrderCtrs(pp, image, loc, order_num)
spirouLOCOR.spirouLOCOR.find_order_centers(pp, image, loc, order_num)
\end{pythonbox}

\begin{pythondocstring}
Find the centre pixels and widths of this order at specific points
along this order="order_num"

specific points are defined by steps (ic_locstepc) away from the
central pixel (ic_cent_col)

:param p: parameter dictionary, ParamDict containing constants
    Must contain at least:
            IC_LOCSTEPC: int, the column separation for fitting orders
            IC_CENT_COL: int, the column number (x-axis) of the central
                         column
            IC_EXT_WINDOW: int, extraction window size (half size)
            IC_IMAGE_GAP: int, the gap index in the selected area
            sigdet: float, the read noise of the image
            IC_LOCSEUIL: float, Normalised amplitude threshold to accept
                         pixels for background calculation
            IC_WIDTHMIN: int, minimum width of order to be accepted
            DRS_DEBUG: int, Whether to run in debug mode
                            0: no debug
                            1: basic debugging on errors
                            2: recipes specific (plots and some code runs)
            DRS_PLOT: bool, Whether to plot (True to plot)

:param image: numpy array (2D), the image

:param loc: parameter dictionary, ParamDict containing data
        Must contain at least:
            ctro: numpy array (2D), storage for the centre positions
                  shape = (number of orders x number of columns (x-axis)

:param order_num: int, the current order to process

:return loc: parameter dictionary, the updated parameter dictionary
        Adds/updates the following:
            ctro: numpy array (2D), storage for the centre positions
                  shape = (number of orders x number of columns (x-axis)
                  updated the values for "order_num"
            sigo: numpy array (2D), storage for the width positions
                  shape = (number of orders x number of columns (x-axis)
                  updated the values for "order_num"
\end{pythondocstring}
\end{minipage}

%----------------------------------------------------------------------------------------
\noindent\begin{minipage}{\textwidth}
\subsection{GetCoeffs}

Defined in \spirouLOCOR\path{.spirouLOCOR.get_loc_coefficients}

\begin{pythonbox}
from SpirouDRS import spirouLOCOR
spirouLOCOR.GetCoeffs(p, hdr=None, loc=None)
spirouLOCOR.spirouLOCOR.get_loc_coefficients(p, hdr=None, loc=None)
\end{pythonbox}

\begin{pythondocstring}
Extracts loco coefficients from parameters keys (uses header="hdr" provided
to get acquisition time or uses p['fitsfilename'] to get acquisition time if
"hdr" is None

:param p: parameter dictionary, ParamDict containing constants
    Must contain at least:
            fitsfilename: string, the full path of for the main raw fits
                          file for a recipe
                          i.e. /data/raw/20170710/filename.fits
            kw_LOCO_NBO: list, keyword store for the number of orders
                         located
            kw_LOCO_DEG_C: list, keyword store for the fit degree for
                           order centers
            kw_LOCO_DEG_W: list, keyword store for the fit degree for
                           order widths
            kw_LOCO_CTR_COEFF: list, keyword store for the coeff center
                               order
            kw_LOCO_FWHM_COEFF: list, keyword store for the coeff width
                                order
            LOC_FILE: string, the suffix for the location calibration
                      database key (usually the fiber type)
                         - read from "loc_file_fpall", if not defined
                           uses p["fiber"]
            fiber: string, the fiber used for this recipe (eg. AB or A or C)
            calibDB: dictionary, the calibration database dictionary
            reduced_dir: string, the reduced data directory
                         (i.e. p['DRS_DATA_REDUC']/p['arg_night_name'])
            log_opt: string, log option, normally the program name

:param hdr: dictionary, header file from FITS rec (opened by spirouFITS)
:param loc: parameter dictionary, ParamDict containing data

:return loc: parameter dictionary, the updated parameter dictionary
        Adds/updates the following:
            number_orders: int, the number of orders in reference spectrum
            nbcoeff_ctr: int, number of coefficients for the centre fit
            nbcoeff_wid: int, number of coefficients for the width fit
            acc: numpy array (2D), the fit coefficients array for
                  the centres fit
                  shape = (number of orders x number of fit coefficients)
            ass: numpy array (2D), the fit coefficients array for
                  the widths fit
\end{pythondocstring}
\end{minipage}

%----------------------------------------------------------------------------------------
\noindent\begin{minipage}{\textwidth}
\subsection{ImageLocSuperimp}

Defined in \spirouLOCOR\path{.spirouLOCOR.image_localization_superposition}

\begin{pythonbox}
from SpirouDRS import spirouLOCOR
spirouLOCOR.ImageLocSuperimp(image, coeffs)
spirouLOCOR.spirouLOCOR.image_localization_superposition(image, coeffs)
\end{pythonbox}

\begin{pythondocstring}
Take an image and superimpose zeros over the positions in the image where
the central fits where found to be

:param image: numpy array (2D), the image
:param coeffs: coefficient array,
               size = (number of orders x number of coefficients in fit)
               output array will be size = (number of orders x dim)
:return newimage: numpy array (2D), the image with super-imposed zero filled
                  fits
\end{pythondocstring}
\end{minipage}

%----------------------------------------------------------------------------------------
\noindent\begin{minipage}{\textwidth}
\subsection{InitialOrderFit}

Defined in \spirouLOCOR\path{.spirouLOCOR.initial_order_fit}

\begin{pythonbox}
from SpirouDRS import spirouLOCOR
spirouLOCOR.InitialOrderFit(pp, loc, mask, onum, rnum, kind, fig=None, frame=None)
spirouLOCOR.spirouLOCOR.initial_order_fit(pp, loc, mask, onum, rnum, kind, fig=None, frame=None)
\end{pythonbox}

\begin{pythondocstring}
Performs a crude initial fit for this order, uses the ctro positions or sigo 
width values found in "FindOrderCtrs" or "find_order_centers" to do the fit

:param p: parameter dictionary, ParamDict containing constants
    Must contain at least:
            log_opt: string, log option, normally the program name
            IC_LOCDFITC: int, order of polynomial to fit for positions
            IC_LOCDFITW: int, order of polynomial to fit for widths
            DRS_PLOT: bool, Whether to plot (True to plot)
            IC_CENT_COL: int, Definition of the central column

:param loc: parameter dictionary, ParamDict containing data
        Must contain at least:
            x: numpy array (1D), the order numbers
            ctro: numpy array (2D), storage for the center positions
                  shape = (number of orders x number of columns (x-axis)
            sigo: numpy array (2D), storage for the width positions
                  shape = (number of orders x number of columns (x-axis)

:param mask: numpy array (1D) of booleans, True where we have non-zero
             widths
:param onum: int, order iteration number (running number over all
             iterations)
:param rnum: int, order number (running number of successful order
             iterations only)
:param kind: string, 'center' or 'fwhm', if 'center' then this fit is for
             the central positions, if 'fwhm' this fit is for the width of
             the orders
:param fig: plt.figure, the figure to plot initial fit on
:param frame: matplotlib axis i.e. plt.subplot(), the axis on which to plot
              the initial fit on (carries the plt.imshow(image))
:return fitdata: dictionary, contains the fit data key value pairs for this
                 initial fit. keys are as follows:

        a = coefficients of the fit from key
        size = 'ic_locdfitc' [for kind='center'] or
             = 'ic_locdftiw' [for kind='fwhm']
        fit = the fity values for the fit (for x = loc['x'])
            where fity = Sum(a[i] * x^i)
        res = the residuals from y - fity
             where y = ctro [kind='center'] or 
                     = sigo [kind='fwhm'])
        abs_res = abs(res)
        rms = the standard deviation of the residuals
        max_ptp = maximum residual value max(res)
        max_ptp_frac = max_ptp / rms  [kind='center']
                     = max(abs_res/y) * 100   [kind='fwhm']
\end{pythondocstring}
\end{minipage}

%----------------------------------------------------------------------------------------
\noindent\begin{minipage}{\textwidth}
\subsection{LocCentralOrderPos}

Defined in \spirouLOCOR\path{.spirouLOCOR.locate_order_center}

\begin{pythonbox}
from SpirouDRS import spirouLOCOR
spirouLOCOR.LocCentralOrderPos(values, threshold, min_width=None)
spirouLOCOR.spirouLOCOR.locate_order_center(values, threshold, min_width=None)
\end{pythonbox}

\begin{pythondocstring}
Takes the values across the oder and finds the order centre by looking for
the start and end of the order (and thus the centre) above threshold

:param values: numpy array (1D) size = number of rows, the pixels in an
                order

:param threshold: float, the threshold above which to find pixels as being
                  part of an order

:param min_width: float, the minimum width for an order to be accepted

:return positions: numpy array (1D), size= number of rows,
                   the pixel positions in cvalues where the centres of each
                   order should be

:return widths:    numpy array (1D), size= number of rows,
                   the pixel positions in cvalues where the centres of each
                   order should be
\end{pythondocstring}
\end{minipage}

%----------------------------------------------------------------------------------------
\noindent\begin{minipage}{\textwidth}
\subsection{MergeCoefficients}

Defined in \spirouLOCOR\path{.spirouLOCOR.merge_coefficients}

\begin{pythonbox}
from SpirouDRS import spirouLOCOR
spirouLOCOR.MergeCoefficients(loc, coeffs, step)
spirouLOCOR.spirouLOCOR.merge_coefficients(loc, coeffs, step)
\end{pythonbox}

\begin{pythondocstring}
Takes a list of coefficients "coeffs" and merges them based on "step"
using the mean of "step" blocks

i.e. shrinks a list of N coefficients to N/2 (if step = 2) where
     indices 0 and 1 are averaged, indices 2 and 3 are averaged etc

:param loc: parameter dictionary, ParamDict containing data
        Must contain at least:
            number_orders: int, the number of orders in reference spectrum

:param coeffs: numpy array (2D), the list of coefficients
               shape = (number of orders x number of fit parameters)

:param step: int, the step between merges
             i.e. total size before = "number_orders"
                  total size after = "number_orders"/step

:return newcoeffs: numpy array (2D), the new list of coefficients
            shape = (number of orders/step x number of fit parameters)
\end{pythondocstring}
\end{minipage}

%----------------------------------------------------------------------------------------
\noindent\begin{minipage}{\textwidth}
\subsection{SigClipOrderFit}

Defined in \spirouLOCOR\path{.spirouLOCOR.sigmaclip_order_fit}

\begin{pythonbox}
from SpirouDRS import spirouLOCOR
spirouLOCOR.SigClipOrderFit(pp, loc, fitdata, mask, onum, rnum, kind)
spirouLOCOR.spirouLOCOR.sigmaclip_order_fit(pp, loc, fitdata, mask, onum, rnum, kind)
\end{pythonbox}

\begin{pythondocstring}
Performs a sigma clip fit for this order, uses the ctro positions or
sigo width values found in "FindOrderCtrs" or "find_order_centers" to do
the fit. Removes the largest residual from the initial fit (or subsequent
sigmaclips) value in x and y and recalculates the fit.

Does this until all the following conditions are NOT met:
       rms > 'ic_max_rms'   [kind='center' or kind='fwhm']
    or max_ptp > 'ic_max_ptp [kind='center']
    or max_ptp_frac > 'ic_ptporms_center'   [kind='center']
    or max_ptp_frac > 'ic_max_ptp_frac'     [kind='fwhm'

:param p: parameter dictionary, ParamDict containing constants
    Must contain at least:
            log_opt: string, log option, normally the program name
            IC_MAX_RMS_CENTER: required when kind="center", float, Maximum
                               rms for sigma-clip order fit (center
                               positions)
            IC_MAX_RMS_FWHM: required when kind="fwhm", float, Maximum
                             rms for sigma-clip order fit (width)
            IC_LOCDFITC: int, order of polynomial to fit for positions
            IC_MAX_PTP_CENTER: required when kind="center", float, Maximum
                               peak-to-peak for sigma-clip order fit
                               (center positions)
            IC_PTPORMS_CENTER: required when kind="center", float, Maximum
                               frac ptp/rms for sigma-clip order fit
                               (center positions)
            IC_LOCDFITW: int, order of polynomial to fit for widths
            IC_MAX_PTP_FRAC_FWHM: required when kind="fwhm", float, Maximum
                                  fractional peak-to-peak for sigma-clip
                                  order fit (width)
            DRS_DEBUG: int, Whether to run in debug mode
                            0: no debug
                            1: basic debugging on errors
                            2: recipes specific (plots and some code runs)
            DRS_PLOT: bool, Whether to plot (True to plot)

:param loc: parameter dictionary, ParamDict containing data
        Must contain at least:
            ctro: numpy array (2D), storage for the center positions
                  shape = (number of orders x number of columns (x-axis)
            sigo: numpy array (2D), storage for the width positions
                  shape = (number of orders x number of columns (x-axis)
            max_rmpts_pos: int, maximum number of removed points in sigma
                           clipping process, for center fits
            max_rmpts_wid: int, maximum number of removed points in sigma
                           clipping process, for width fits

\end{pythondocstring}
\end{minipage}

\noindent\begin{minipage}{\textwidth}                        
\begin{pythondocstring}
sigmaclip_order_fit (continued)

:param fitdata: dictionary, contains the fit data key value pairs for this
                 initial fit. keys are as follows:

        a = coefficients of the fit from key
        size = 'ic_locdfitc' [for kind='centre'] or
             = 'ic_locdftiw' [for kind='fwhm']
        fit = the fity values for the fit (for x = loc['x'])
            where fity = Sum(a[i] * x^i)
        res = the residuals from y - fity
             where y = ctro [kind='center'] or
                     = sigo [kind='fwhm'])
        abs_res = abs(res)
        rms = the standard deviation of the residuals
        max_ptp = maximum residual value max(res)
        max_ptp_frac = max_ptp / rms  [kind='center']
                     = max(abs_res/y) * 100   [kind='fwhm']

:param mask: numpy array (1D) of booleans, True where we have non-zero
             widths
:param onum: int, order iteration number (running number over all
             iterations)
:param rnum: int, order number (running number of successful order
             iterations only)
:param kind: string, 'center' or 'fwhm', if 'center' then this fit is for
             the central p

:return fitdata: dictionary, contains the fit data key value pairs for this
                 initial fit. keys are as follows:

        a = coefficients of the fit from key
        size = 'ic_locdfitc' [for kind='center'] or
             = 'ic_locdftiw' [for kind='fwhm']
        fit = the fity values for the fit (for x = loc['x'])
            where fity = Sum(a[i] * x^i)
        res = the residuals from y - fity
             where y = ctro [kind='center'] or
                     = sigo [kind='fwhm'])
        abs_res = abs(res)
        rms = the standard deviation of the residuals
        max_ptp = maximum residual value max(res)
        max_ptp_frac = max_ptp / rms  [kind='center']
                     = max(abs_res/y) * 100   [kind='fwhm']
\end{pythondocstring}
\end{minipage}

%----------------------------------------------------------------------------------------

%%%%%%%%%%%%%%%%%%%%%%%%%%%%%%%%%%%%%%%%%%%%%%%%%%%%%%%%
%%
\clearpage
\newpage
\noindent\begin{minipage}{\textwidth}
\section{The spirouPOLAR module}
\label{ch:the_module:spirouPOLAR}
%%
%%%%%%%%%%%%%%%%%%%%%%%%%%%%%%%%%%%%%%%%%%%%%%%%%%%%%%%%

%----------------------------------------------------------------------------------------
\subsection{SortPolarFiles}

Defined in \spirouPOLAR\path{.spirouPOLAR.sort_polar_files}

\begin{pythonbox}
from SpirouDRS import spirou
spirouPOLAR.SortPolarFiles(p, polardict)
spirouPOLAR.spirouPOLAR.sort_polar_files(p, polardict)
\end{pythonbox}

\begin{pythondocstring}
Doc string missing!
\end{pythondocstring}
\end{minipage}


%----------------------------------------------------------------------------------------
\noindent\begin{minipage}{\textwidth}
\subsection{LoadPolarData}

Defined in \spirouPOLAR\path{.spirouPOLAR.load_data}

\begin{pythonbox}
from SpirouDRS import spirou
spirouPOLAR.LoadPolarData(p, polardict, loc)
spirouPOLAR.spirouPOLAR.load_data(p, polardict, loc)
\end{pythonbox}

\begin{pythondocstring}
Doc string missing!
\end{pythondocstring}
\end{minipage}


%----------------------------------------------------------------------------------------

\noindent\begin{minipage}{\textwidth}
\subsection{CalculatePolarimetry}

Defined in \spirouPOLAR\path{.spirouPOLAR.calculate_polarimetry_wrapper}

\begin{pythonbox}
from SpirouDRS import spirou
spirouPOLAR.CalculatePolarimetry(p, loc)
spirouPOLAR.spirouPOLAR.calculate_polarimetry_wrapper(p, loc)
\end{pythonbox}

\begin{pythondocstring}
Doc string missing!
\end{pythondocstring}
\end{minipage}


%----------------------------------------------------------------------------------------

\noindent\begin{minipage}{\textwidth}
\subsection{CalculateContinuum}

Defined in \spirouPOLAR\path{.spirouPOLAR.calculate_continuum}

\begin{pythonbox}
from SpirouDRS import spirou
spirouPOLAR.CalculateContinuum(p, loc, in_wavelength=True)
spirouPOLAR.spirouPOLAR.calculate_continuum(p, loc, in_wavelength=True)
\end{pythonbox}

\begin{pythondocstring}
Doc string missing!
\end{pythondocstring}
\end{minipage}


%----------------------------------------------------------------------------------------

\noindent\begin{minipage}{\textwidth}
\subsection{CalculateStokesI}

Defined in \spirouPOLAR\path{.spirouPOLAR.calculate_stokes_I}

\begin{pythonbox}
from SpirouDRS import spirou
spirouPOLAR.CalculateStokesI(p, loc)
spirouPOLAR.spirouPOLAR.calculate_stokes_I(p, loc)
\end{pythonbox}

\begin{pythondocstring}
Doc string missing!
\end{pythondocstring}
\end{minipage}


%----------------------------------------------------------------------------------------



%%%%%%%%%%%%%%%%%%%%%%%%%%%%%%%%%%%%%%%%%%%%%%%%%%%%%%%%
%%
\clearpage
\newpage
\begin{minipage}{\textwidth}
\section{The  module}
\label{ch:the_module:}
%%
%%%%%%%%%%%%%%%%%%%%%%%%%%%%%%%%%%%%%%%%%%%%%%%%%%%%%%%%

%----------------------------------------------------------------------------------------
\subsection{}

Defined in \spirouRV\path{.}

\begin{pythonbox}
from SpirouDRS import spirouRV
spirouRV.
spirouRV.
\end{pythonbox}

\begin{pythondocstring}

\end{pythondocstring}
\end{minipage}

%----------------------------------------------------------------------------------------
\begin{minipage}{\textwidth}
\subsection{}

Defined in \spirouRV\path{.}

\begin{pythonbox}
from SpirouDRS import spirouRV
spirouRV.
spirouRV.
\end{pythonbox}

\begin{pythondocstring}

\end{pythondocstring}
\end{minipage}

%----------------------------------------------------------------------------------------
\begin{minipage}{\textwidth}
\subsection{}

Defined in \spirouRV\path{.}

\begin{pythonbox}
from SpirouDRS import spirouRV
spirouRV.
spirouRV.
\end{pythonbox}

\begin{pythondocstring}

\end{pythondocstring}
\end{minipage}

%----------------------------------------------------------------------------------------
\begin{minipage}{\textwidth}
\subsection{}

Defined in \spirouRV\path{.}

\begin{pythonbox}
from SpirouDRS import spirouRV
spirouRV.
spirouRV.
\end{pythonbox}

\begin{pythondocstring}

\end{pythondocstring}
\end{minipage}

%----------------------------------------------------------------------------------------
\begin{minipage}{\textwidth}
\subsection{}

Defined in \spirouRV\path{.}

\begin{pythonbox}
from SpirouDRS import spirouRV
spirouRV.
spirouRV.
\end{pythonbox}

\begin{pythondocstring}

\end{pythondocstring}
\end{minipage}

%----------------------------------------------------------------------------------------
\begin{minipage}{\textwidth}
\subsection{}

Defined in \spirouRV\path{.}

\begin{pythonbox}
from SpirouDRS import spirouRV
spirouRV.
spirouRV.
\end{pythonbox}

\begin{pythondocstring}

\end{pythondocstring}
\end{minipage}

%----------------------------------------------------------------------------------------
\begin{minipage}{\textwidth}
\subsection{}

Defined in \spirouRV\path{.}

\begin{pythonbox}
from SpirouDRS import spirouRV
spirouRV.
spirouRV.
\end{pythonbox}

\begin{pythondocstring}

\end{pythondocstring}
\end{minipage}

%----------------------------------------------------------------------------------------
\begin{minipage}{\textwidth}
\subsection{}

Defined in \spirouRV\path{.}

\begin{pythonbox}
from SpirouDRS import spirouRV
spirouRV.
spirouRV.
\end{pythonbox}

\begin{pythondocstring}

\end{pythondocstring}
\end{minipage}

%----------------------------------------------------------------------------------------
\begin{minipage}{\textwidth}
\subsection{}

Defined in \spirouRV\path{.}

\begin{pythonbox}
from SpirouDRS import spirouRV
spirouRV.
spirouRV.
\end{pythonbox}

\begin{pythondocstring}

\end{pythondocstring}
\end{minipage}

%----------------------------------------------------------------------------------------
\begin{minipage}{\textwidth}
\subsection{}

Defined in \spirouRV\path{.}

\begin{pythonbox}
from SpirouDRS import spirouRV
spirouRV.
spirouRV.
\end{pythonbox}

\begin{pythondocstring}

\end{pythondocstring}
\end{minipage}

%----------------------------------------------------------------------------------------
\begin{minipage}{\textwidth}
\subsection{}

Defined in \spirouRV\path{.}

\begin{pythonbox}
from SpirouDRS import spirouRV
spirouRV.
spirouRV.
\end{pythonbox}

\begin{pythondocstring}

\end{pythondocstring}
\end{minipage}

%----------------------------------------------------------------------------------------
\begin{minipage}{\textwidth}
\subsection{}

Defined in \spirouRV\path{.}

\begin{pythonbox}
from SpirouDRS import spirouRV
spirouRV.
spirouRV.
\end{pythonbox}

\begin{pythondocstring}

\end{pythondocstring}
\end{minipage}

%----------------------------------------------------------------------------------------
\begin{minipage}{\textwidth}
\subsection{}

Defined in \spirouRV\path{.}

\begin{pythonbox}
from SpirouDRS import spirouRV
spirouRV.
spirouRV.
\end{pythonbox}

\begin{pythondocstring}

\end{pythondocstring}
\end{minipage}

%----------------------------------------------------------------------------------------
\begin{minipage}{\textwidth}
\subsection{}

Defined in \spirouRV\path{.}

\begin{pythonbox}
from SpirouDRS import spirouRV
spirouRV.
spirouRV.
\end{pythonbox}

\begin{pythondocstring}

\end{pythondocstring}
\end{minipage}

%----------------------------------------------------------------------------------------



%%%%%%%%%%%%%%%%%%%%%%%%%%%%%%%%%%%%%%%%%%%%%%%%%%%%%%%%
%%
\clearpage
\newpage
\begin{minipage}{\textwidth}
\section{The spirouStartup module}
\label{ch:the_module:spirouStartup}
%%
%%%%%%%%%%%%%%%%%%%%%%%%%%%%%%%%%%%%%%%%%%%%%%%%%%%%%%%%

%----------------------------------------------------------------------------------------
\subsection{Begin}

Defined in \spirouStartup\path{.spirouStartup.run_begin}

\begin{pythonbox}
from SpirouDRS import spirouStartup
spirouStartup.Begin()
spirouStartup.spirouStartup.run_begin()
\end{pythonbox}

\begin{pythondocstring}
Begin DRS - Must be run at start of every recipe
- loads the parameters from the primary configuration file, displays
  title, checks priamry constants and displays initial parameterization

:return cparams: parameter dictionary, ParamDict constants from primary
                 configuration file
        Adds the following:
            all constants in primary configuration file
            DRS_NAME: string, the name of the DRS
            DRS_VERSION: string, the version of the DRS
\end{pythondocstring}
\end{minipage}

%----------------------------------------------------------------------------------------
\begin{minipage}{\textwidth}
\subsection{GetCustomFromRuntime}

Defined in \spirouStartup\path{.spirouStartup.get_custom_from_run_time_args}

\begin{pythonbox}
from SpirouDRS import spirouStartup
spirouStartup.GetCustomFromRuntime(positions=None, types=None, names=None,
                                   required=None, calls=None, cprior=None,
                                   lognames=None)
spirouStartup.spirouStartup.get_custom_from_run_time_args(positions=None, types=None, names=None,
                                  required=None, calls=None, cprior=None,
                                  lognames=None)
\end{pythonbox}

\begin{pythondocstring}
Extract custom arguments from defined positions in sys.argv (defined at
run time)

:param positions: list of integers, the positions of the arguments
                  (i.e. first argument is 0)

:param types: list of python types, the type (i.e. int, float) for each
              argument
:param names: list of strings, the names of each argument (to access in
              parameter dictionary once extracted)

:param required: list of bools or None, states whether the program
                 should exit if runtime argument not found

:param calls: list of objects or None, if define these are the values that
              come from a function call (overwrite command line arguments)

:param lognames: list of strings, the names displayed in the log (on error)
                 theses should be similar to "names" but in a form the
                 user can easily understand for each variable

:return values: dictionary, if run time arguments are correct python type
                the name-value pairs are returned
\end{pythondocstring}
\end{minipage}

%----------------------------------------------------------------------------------------
\begin{minipage}{\textwidth}
\subsection{GetFile}

Defined in \spirouStartup\path{.spirouStartup.get_file}

\begin{pythonbox}
from SpirouDRS import spirouStartup
spirouStartup.GetFile(p, path, name=None, prefix=None, kind=None)
spirouStartup.spirouStartup.get_file(p, path, name=None, prefix=None, kind=None)
\end{pythonbox}

\begin{pythondocstring}
Get full file path and check the path and file exist

:param p: parameter dictionary, ParamDict containing constants
    Must contain at least:
            log_opt: string, log option, normally the program name
            program: string, the recipe/way the script was called
                     i.e. from sys.argv[0]

:param path: string, either the directory to the folder (if name is None) or
             the full path to the file
:param name: string or None, the name of the file, if None name is assumed
             to be in path
:param prefix: string or None, if not None this substring must be in the
               filename
:param kind: string or None, the type of file (for logging)

:return location: string, the full file path of the file
\end{pythondocstring}
\end{minipage}

%----------------------------------------------------------------------------------------

\begin{minipage}{\textwidth}
\subsection{GetFiberType}

Defined in \spirouStartup\path{.spirouStartup.get_fiber_type}

\begin{pythonbox}
from SpirouDRS import spirouStartup
spirouStartup.GetFiberType(p, filename, fibertypes=None)
spirouStartup.spirouStartup.get_fiber_type(p, filename, fibertypes=None)
\end{pythonbox}

\begin{pythondocstring}
Get fiber types and search for a valid fiber type in filename

:param p: parameter dictionary, ParamDict containing constants
    Must contain at least:
            FIBER_TYPES: list of strings, the types of fiber available
                         (i.e. ['AB', 'A', 'B', 'C'])
            log_opt: string, log option, normally the program name

:param filename: string, the filename to search for fiber types in
:param fibertypes: list of strings, the fiber types to search for

:return fiber: string, the fiber found (exits via WLOG if no fiber found)
\end{pythondocstring}
\end{minipage}

%----------------------------------------------------------------------------------------

\begin{minipage}{\textwidth}
\subsection{LoadArguments}

Defined in \spirouStartup\path{.spirouStartup.load_arguments}

\begin{pythonbox}
from SpirouDRS import spirouStartup
spirouStartup.LoadArguments(cparams, night_name=None, files=None, customargs=None)
spirouStartup.spirouStartup.load_arguments(cparams, night_name=None, files=None, customargs=None)
\end{pythonbox}

\begin{pythondocstring}
Deal with loading run time arguments:

1) display help file (if requested and exists)
2) loads run time arguments (and custom arguments, see below)
3) loads other config files

:param cparams: parameter dictionary, ParamDict containing constants
    Must contain at least:
            arg_night_name: string, the folder within data raw directory
                            containing files (also reduced directory) i.e.
                            /data/raw/20170710 would be "20170710"
            arg_file_names: list, list of files taken from the command line
                            (or call to recipe function) must have at least
                            one string filename in the list

:param night_name: string or None, the name of the directory in DRS_DATA_RAW
                   to find the files in

                   if None (undefined) uses the first argument in command
                   line (i.e. sys.argv[1])

                   if defined overwrites call from
                   command line (i.e. overwrites sys.argv)

                   stored in p['arg_night_name']

:param files: list of strings or None, the files to use for this program

              if None (undefined) uses the second and all other arguments in
              the command line (i.e. sys.argv[2:])

              if defined overwrites call from command line

              stored in p['arg_file_names']

:param customargs: None or list of strings, if list of strings then instead
                   of getting the standard runtime arguments

       i.e. in form:

            program.py rawdirectory arg_file_names[0] arg_file_names[1]...

       loads all arguments into customargs

       i.e. if customargs = ['rawdir', 'filename', 'a', 'b', 'c']
       expects command line arguments to be:

            program.py rawdir filename a b c

:return p: dictionary, parameter dictionary
\end{pythondocstring}
\end{minipage}

%----------------------------------------------------------------------------------------

\begin{minipage}{\textwidth}
\subsection{InitialFileSetup}

Defined in \spirouStartup\path{.spirouStartup.initial_file_setup}

\begin{pythonbox}
from SpirouDRS import spirouStartup
spirouStartup.InitialFileSetup(p, kind=None, prefixes=None, add_to_p=None,
                               calibdb=False)
spirouStartup.spirouStartup.initial_file_setup(p, kind=None, prefixes=None, add_to_p=None,
                                               calibdb=False)
\end{pythonbox}

\begin{pythondocstring}
Run start up code (based on program and parameters defined in p before)

:param p: parameter dictionary, ParamDict containing constants
    Must contain at least:
            log_opt: string, log option, normally the program name
            fitsfilename: string, the full path of for the main raw fits
                          file for a recipe
                          i.e. /data/raw/20170710/filename.fits
            program: string, the recipe/way the script was called
                     i.e. from sys.argv[0]
            reduced_dir: string, the reduced data directory
                         (i.e. p['DRS_DATA_REDUC']/p['arg_night_name'])
            DRS_DATA_REDUC: string, the directory that the reduced data
                            should be saved to/read from
            DRS_CALIB_DB: string, the directory that the calibration
                          files should be saved to/read from

:param kind: string, description of program we are running (i.e. dark)

:param prefixes: list of strings, prefixes to look for in file name
:param prefixes: list of strings, prefixes to look for in file name
                 will exit code if none of the prefixes are found
                 (prefix = None if no prefixes are needed to be found)

:param add_to_p: dictionary structure:

        add_to_p[prefix1] = dict(key1=value1, key2=value2)
        add_to_p[prefix2] = dict(key3=value3, key4=value4)

        where prefix1 and prefix2 are the strings in "prefixes"

        This will add the sub dictionarys to the main parameter dictionary
        based on which prefix is found

        i.e. if prefix1 is found key "value3" and "value4" above are added
        (with "key3" and "key4") to the parameter dictionary p

:param calibdb: bool, if True calibDB folder and files are required and
                program will log and exit if they are not found
                if False, program will create calibDB folder

:return p: parameter dictionary, the updated parameter dictionary
        Adds the following:
            calibDB: dictionary, the calibration database dictionary
            prefixes from add_to_p (see spirouStartup.deal_with_prefixes)
\end{pythondocstring}
\end{minipage}

%----------------------------------------------------------------------------------------




%%%%%%%%%%%%%%%%%%%%%%%%%%%%%%%%%%%%%%%%%%%%%%%%%%%%%%%%
%%
\clearpage
\newpage
\noindent\begin{minipage}{\textwidth}
\section{The spirouTHORCA module}
\label{ch:the_module:spirouTHORCA}
%%
%%%%%%%%%%%%%%%%%%%%%%%%%%%%%%%%%%%%%%%%%%%%%%%%%%%%%%%%

%----------------------------------------------------------------------------------------
\subsection{CalcInstrumentDrift}

Defined in \spirouTHORCA\path{.spirouWAVE.calculate_instrument_drift}

\begin{pythonbox}
from SpirouDRS import spirouTHORCA
spirouTHORCA.CalcInstrumentDrift
spirouTHORCA.spirouWAVE.calculate_instrument_drift

\end{pythonbox}

\begin{pythondocstring}
Calculate the instrumental drift between the reference file and the current
file

:param p: parameter dictionary, ParamDict containing constants
    Must contain at least:
            log_opt: string, log option, normally the program name
            DRS_PLOT: bool, if True do plots else do not
            FIBER: string, the fiber type (i.e. AB or A or B or C)
            IC_HC_N_ORD_FINAL: int, the order to which the solution is
                               calculated
            IC_WAVE_IDRIFT_NOISE: float, the noise used in drift calculation
            IC_WAVE_IDRIFT_BOXSIZE: int, the size around a saturated pixel
                                    to count as bad
            IC_WAVE_IDRIFT_MAXFLUX: int, the maximum flux for a good
                                    (unsaturated) pixel
            IC_WAVE_IDRIFT_RV_CUT: float, the rv cut above which rv from
                                   orders are not used
            QC_WAVE_IDRIFT_NBORDEROUT: int, the maximum number of orders to
                                       remove from RV calculation
            QC_WAVE_IDRIFT_RV_MAX: float, the maximum allowed drift (in m/s)

:param loc: parameter dictionary, ParamDict containing data
    Must contain at least:
        HCDATA: numpy array (2D), the image data
        HCHDR: dictionary, the header dictionary for HCDATA

:return loc: parameter dictionary, the updated parameter dictionary
        Adds/updates the following:
            DRIFT_REF: string, the filename of the reference data
            DRIFT_RV: float, the mean instrumental RV
            DRIFT_NBCOS: int, the number of cosmic pixels found
            DRIFT_RFLUX: float, the mean flux ratio
            DRIFT_NBORDKILL: int, the number of orders removed
            DRIFT_NOISE: flat, the weighted mean uncertainty on the RV
\end{pythondocstring}
\end{minipage}


%----------------------------------------------------------------------------------------
\noindent\begin{minipage}{\textwidth}
\subsection{CalcLittrowSolution}

Defined in \spirouTHORCA\path{.spirouTHORCA.calculate_littrow_sol}

\begin{pythonbox}
from SpirouDRS import spirouTHORCA
spirouTHORCA.CalcLittrowSolution
spirouTHORCA.spirouTHORCA.calculate_littrow_sol

\end{pythonbox}

\begin{pythondocstring}
Calculate the Littrow solution for this iteration for a set of cut points

Uses ALL_LINES_i  where i = iteration to calculate the littrow solutions
for defiend cut points (given a cut_step and fit_deg of
IC_LITTROW_CUT_STEP_i and IC_LITTROW_FIT_DEG_i where i = iteration)

:param p: parameter dictionary, ParamDict containing constants
    Must contain at least:
        LOG_OPT: string, log option, normally the program name
        FIBER: string, the fiber type (i.e. AB or A or B or C)
        IC_LITTROW_REMOVE_ORDERS: list of ints, if not empty removes these
                                  orders from influencing the fit 
        IC_LITTROW_ORDER_INIT: int, defines the first order to for the fit
                               solution
        IC_HC_N_ORD_START: int, defines first order HC solution was
                           calculated from
        IC_HC_N_ORD_FINAL: int, defines last order HC solution was
                           calculated to
        IC_LITTROW_CUT_STEP_i: int, defines the step to use between 
                               cut points 
        IC_LITTROW_FIT_DEG_i: int, defines the polynomial fit degree

        where i = iteration

:param loc: parameter dictionary, ParamDict containing data
    Must contain at least:
        ECHELLE_ORDERS: numpy array (1D), the echelle order numbers
        HCDATA: numpy array (2D), the image data (used for shape)
        ALL_LINES_i: list of numpy arrays, length = number of orders
                     each numpy array contains gaussian parameters
                     for each found line in that order

        where i = iteration

:param ll: numpy array (1D), the initial guess wavelengths for each line
:param iteration: int, the iteration number (used so we can store multiple
                  calculations in loc, defines "i" in input and outputs
                  from p and loc
:param log: bool, if True will print a final log message on completion with
            some stats

:return loc: parameter dictionary, the updated parameter dictionary
        Adds/updates the following:
            X_CUT_POINTS_i: numpy array (1D), the x pixel cut points
            LITTROW_MEAN_i: list, the mean position of each cut point
            LITTROW_SIG_i: list, the mean FWHM of each cut point
            LITTROW_MINDEV_i: list, the minimum deviation of each cut point
            LITTROW_MAXDEV_i: list, the maximum deviation of each cut point
            LITTROW_PARAM_i: list of numpy arrays, the gaussian fit
                             coefficients of each cut point
            LITTROW_XX_i: list, the order positions of each cut point
            LITTROW_YY_i: list, the residual fit of each cut point

            where i = iteration

ALL_LINES_i definition:
    ALL_LINES_i[row] = [gparams1, gparams2, ..., gparamsN]

                where:
                    gparams[0] = output wavelengths
                    gparams[1] = output sigma (gauss fit width)
                    gparams[2] = output amplitude (gauss fit)
                    gparams[3] = difference in input/output wavelength
                    gparams[4] = input amplitudes
                    gparams[5] = output pixel positions
                    gparams[6] = output pixel sigma width
                                      (gauss fit width in pixels)
                    gparams[7] = output weights for the pixel position
\end{pythondocstring}
\end{minipage}


%----------------------------------------------------------------------------------------
\noindent\begin{minipage}{\textwidth}
\subsection{DetectBadLines}

Defined in \spirouTHORCA\path{.spirouTHORCA.detect_bad_lines}

\begin{pythonbox}
from SpirouDRS import spirouTHORCA
spirouTHORCA.DetectBadLines(p, loc, key=None, iteration=0)
spirouTHORCA.spirouTHORCA.detect_bad_lines(p, loc, key=None, iteration=0)
\end{pythonbox}

\begin{pythondocstring}
Detect and filter out the bad lines for this iteration

:param p: parameter dictionary, ParamDict containing constants
    Must contain at least:
        IC_MAX_ERRW_ONFIT: float, the maximum error on the first guess
                           lines
        IC_MAX_SIGLL_CAL_LINES: float, the maximum sigma fit of the guessed
                                lines
        IC_MAX_AMPL_LINE: float, the maximum amplitude the guessed lines
                          can have
        IC_HC_T_ORDER_START: int, the echelle number of the first
                             extracted order

:param loc: parameter dictionary, ParamDict containing data
    Must contain at least:
        ECHELLE_ORDERS: numpy array (1D), the echelle order numbers
        ALL_LINES_i or "key": if key is None use ALL_LINES_i
                              where i = iteration
                     list of numpy arrays, length = number of orders
                     each numpy array contains gaussian parameters
                     for each found line in that order

:param key: string or None, if string should exist in "loc" and replaces
            "ALL_LINES_i" where i = iteration, and should point to a list
            of numpy arrays, length = number of orders, each numpy array
            contains gaussian parameters for each found line in that order
:param iteration: int, the iteration number (used so we can store multiple
                  calculations in loc, defines "i" in input and outputs
                  from p and loc

:return loc: parameter dictionary, the updated parameter dictionary
    Adds/updates the following:
        ALL_LINES_i or "key": list of numpy arrays, ALL_LINES_i or key input
                              with the bad lines filtered out

ALL_LINES_i definition:
    ALL_LINES_i[row] = [gparams1, gparams2, ..., gparamsN]

                where:
                    gparams[0] = output wavelengths
                    gparams[1] = output sigma (gauss fit width)
                    gparams[2] = output amplitude (gauss fit)
                    gparams[3] = difference in input/output wavelength
                    gparams[4] = input amplitudes
                    gparams[5] = output pixel positions
                    gparams[6] = output pixel sigma width
                                      (gauss fit width in pixels)
                    gparams[7] = output weights for the pixel position
\end{pythondocstring}
\end{minipage}

%----------------------------------------------------------------------------------------
\noindent\begin{minipage}{\textwidth}
\subsection{ExtrapolateLittrowSolution}

Defined in \spirouTHORCA\path{.spirouTHORCA.extrapolate_littrow_sol}

\begin{pythonbox}
from SpirouDRS import spirouTHORCA
spirouTHORCA.ExtrapolateLittrowSolution(p, loc, ll, iteration=0)
spirouTHORCA.spirouTHORCA.extrapolate_littrow_sol(p, loc, ll, iteration=0)
\end{pythonbox}

\begin{pythondocstring}
Extrapolate and fit the Littrow solution at defined points and return
the wavelengths, solutions, and cofficients of the littorw fits

:param p: parameter dictionary, ParamDict containing constants
    Must contain at least:
        IC_LITTROW_ORDER_FIT_DEG: int, defines the polynomial fit degree
        IC_HC_T_ORDER_START
        IC_LITTROW_ORDER_INIT int, defines the first order to for the fit
                               solution
:param loc: parameter dictionary, ParamDict containing data
    Must contain at least:
        HCDATA: numpy array (2D), the image data (used for shape)
        LITTROW_PARAM_i: list of numpy arrays, the gaussian fit
                         coefficients of each cut point
        X_CUT_POINTS_i: numpy array (1D), the x pixel cut points
    
        where i = iteration
        
:param ll: numpy array (1D), the initial guess wavelengths for each line
:param iteration: int, the iteration number (used so we can store multiple
                  calculations in loc, defines "i" in input and outputs
                  from p and loc
                  
:return loc: parameter dictionary, the updated parameter dictionary
        Adds/updates the following:
            LITTROW_EXTRAP_i: numpy array (2D), 
                              size=([no. orders] by [no. cut points])
                              the wavelength values at each cut point for
                              each order
            LITTROW_EXTRAP_SOL_i: numpy array (2D), 
                              size=([no. orders] by [no. cut points])
                              the wavelength solution at each cut point for
                              each order
            LITTROW_EXTRAP_PARAM_i: numy array (2D),
                              size=([no. orders] by [the fit degree +1])
                              the coefficients of the fits for each cut 
                              point for each order
            
            where i = iteration
\end{pythondocstring}
\end{minipage}


%----------------------------------------------------------------------------------------


\noindent\begin{minipage}{\textwidth}
\subsection{FirstGuessSolution}

Defined in \spirouTHORCA\path{.spirouTHORCA.first_guess_at_wave_solution}

\begin{pythonbox}
from SpirouDRS import spirouTHORCA
spirouTHORCA.FirstGuessSolution(p, loc, mode=0)
spirouTHORCA.spirouTHORCA.first_guess_at_wave_solution(p, loc, mode=0)
\end{pythonbox}

\begin{pythondocstring}
First guess at wave solution, consistency check, using the wavelength
solutions line list

:param p: parameter dictionary, ParamDict containing constants
    Must contain at least:
        IC_HC_N_ORD_START: int, defines first order solution is calculated
        IC_HC_N_ORD_FINAL: int, defines last order solution is calculated
                            from
        IC_HC_T_ORDER_START: int, defines the echelle order of
                            the first e2ds order
        log_opt: string, log option, normally the program name
        fiber: string, the fiber number

:param loc: parameter dictionary, ParamDict containing data
    Must contain at least:
        ECHELLE_ORDERS: numpy array (1D), the echelle order numbers

:param mode: string, if mode="new" uses python to work out gaussian fit
                     if mode="old" uses FORTRAN (testing only) requires
                     compiling of FORTRAN fitgaus into spirouTHORCA dir

:return loc: parameter dictionary, the updated parameter dictionary
        Adds/updates the following:
            ECHELLE_ORDERS: numpy array, the echelle orders to fit
            LL_INIT: numpy array, the initial guess at the line list
            LL_LINE: numpy array, the line list wavelengths from file
            AMPL_LINE: numpy array, the line list amplitudes from file
            ALL_LINES: list of numpy arrays, length = number of orders
                        each numpy array contains gaussian parameters
                        for each found line in that order

ALL_LINES_i definition:
    ALL_LINES_i[row] = [gparams1, gparams2, ..., gparamsN]

                where:
                    gparams[0] = output wavelengths
                    gparams[1] = output sigma (gauss fit width)
                    gparams[2] = output amplitude (gauss fit)
                    gparams[3] = difference in input/output wavelength
                    gparams[4] = input amplitudes
                    gparams[5] = output pixel positions
                    gparams[6] = output pixel sigma width
                                      (gauss fit width in pixels)
                    gparams[7] = output weights for the pixel position
\end{pythondocstring}
\end{minipage}

%----------------------------------------------------------------------------------------
\noindent\begin{minipage}{\textwidth}
\subsection{Fit1DSolution}

Defined in \spirouTHORCA\path{.spirouTHORCA.fit_1d_solution}

\begin{pythonbox}
from SpirouDRS import spirouTHORCA
spirouTHORCA.Fit1DSolution(p, loc, ll, iteration=0)
spirouTHORCA.spirouTHORCA.fit_1d_solution(p, loc, ll, iteration=0)
\end{pythonbox}

\begin{pythondocstring}
Fits the 1D solution between wavelength and pixel postion and inverts it
into a fit between pixel position and wavelength

:param p: parameter dictionary, ParamDict containing constants
    Must contain at least:
        LOG_OPT: string, log option, normally the program name
        FIBER: string, the fiber type (i.e. AB or A or B or C)

:param loc: parameter dictionary, ParamDict containing data
    Must contain at least:
        ECHELLE_ORDERS: numpy array (1D), the echelle order numbers
        HCDATA: numpy array (2D), the image data (used for shape)
        ALL_LINES_i: list of numpy arrays, length = number of orders
                     each numpy array contains gaussian parameters
                     for each found line in that order
        IC_ERRX_MIN: float, the minimum instrumental error
        IC_LL_DEGR_FIT: int, the wavelength fit polynomial order
        IC_MAX_LLFIT_RMS: float, the max rms for the wavelength
                          sigma-clip fit
        IC_HC_T_ORDER_START: int, defines the echelle order of
                            the first e2ds order

:param ll: numpy array (1D), the initial guess wavelengths for each line
:param iteration: int, the iteration number (used so we can store multiple
                  calculations in loc, defines "i" in input and outputs
                  from p and loc

:return loc: parameter dictionary, the updated parameter dictionary
        Adds/updates the following:
            X_MEAN_i: float, the mean residual from the fit [km/s]
            X_VAR_i: float, the rms residual from the fit [km/s]
            X_ITER_i: numpy array(1D), the last line central position, FWHM
            and the number of lines in each order
            X_PARAM_i: numpy array (1D), the coefficients of the fit to x
                       pixel position as a function of wavelength
            X_DETAILS_i: list, [lines, xfit, cfit, weight] where
                         lines= original wavelength-centers used for the fit
                         xfit= original pixel-centers used for the fit
                         cfit= fitted pixel-centers using fit coefficients
                         weight=the line weights used
            SCALE_i: list, the convertion for each line into wavelength
            LL_MEAN_i: float, the mean residual after inversion [km/s]
            LL_VAR_i: float, the rms residual after inversion [km/s]
            LL_PARAM_i: numpy array (1D), the cofficients of the inverted
                        fit (i.e. wavelength as a function of x pixel
                        position)
            LL_DETAILS_i: numpy array (1D), the [nres, wei] where
                          nres = normalised residuals in km/s
                          wei = the line weights
            LL_OUT_i: numpy array (2D), the output wavelengths for each
                      pixel and each order (in the shape of original image)
            DLL_OUT_i: numpy array (2D), the output delta wavelengths for
                       each pixel and each order (in the shape of original
                       image)

            where i = iteration

ALL_LINES_i definition:
    ALL_LINES_i[row] = [gparams1, gparams2, ..., gparamsN]

                where:
                    gparams[0] = output wavelengths
                    gparams[1] = output sigma (gauss fit width)
                    gparams[2] = output amplitude (gauss fit)
                    gparams[3] = difference in input/output wavelength
                    gparams[4] = input amplitudes
                    gparams[5] = output pixel positions
                    gparams[6] = output pixel sigma width
                                      (gauss fit width in pixels)
                    gparams[7] = output weights for the pixel position
\end{pythondocstring}
\end{minipage}

%----------------------------------------------------------------------------------------
\noindent\begin{minipage}{\textwidth}
\subsection{FPWaveSolution}

Defined in \spirouTHORCA\path{.spirouWAVE.fp_wavelength_sol}

\begin{pythonbox}
from SpirouDRS import spirouTHORCA
spirouTHORCA.FPWaveSolution(p, loc, mode='new')
spirouTHORCA.spirouWAVE.fp_wavelength_sol(p, loc, mode='new')
\end{pythonbox}

\begin{pythondocstring}
Derives the FP line wavelengths from the first solution
Follows the Bauer et al 2015 procedure

:param p: parameter dictionary, ParamDict containing constants
    Must contain at least:
        IC_FP_N_ORD_START: int, defines first order FP solution is
                           calculated from

        IC_FP_N_ORD_FINAL: int, defines last order FP solution is
                           calculated to

        IC_HC_N_ORD_START_2: int, defines first order HC solution was
                             calculated from

        IC_HC_N_ORD_FINAL_2: int, defines last order HC solution was
                             calculated to

        IC_FP_THRESHOLD: float, defines threshold for detecting FP lines

        IC_FP_SIZE: int, defines size of region where each line is fitted

        IC_FP_DOPD0: float, initial value of FP effective cavity width

        IC_FP_FIT_DEGREE: int, degree of polynomial fit

        log_opt: string, log option, normally the program name

:param loc: parameter dictionary, ParamDict containing data
    Must contain at least:
        FP_DATA: the FP e2ds data
        LITTROW_EXTRAP_SOL_1: the wavelength solution derived from the HC
                              and Littrow-constrained
        ALL_LINES: list of numpy arrays, length = number of orders
                   each numpy array contains gaussian parameters
                   for each found line in that order
        BLAZE: numpy array (2D), the blaze data

:return loc: parameter dictionary, the updated parameter dictionary
        Adds/updates the following:
            FP_LL_POS: numpy array, the initial wavelengths of the FP lines
            FP_XX_POS: numpy array, the pixel positions of the FP lines
            FP_M: numpy array, the FP line numbers
            FP_DOPD_T: numpy array, the measured cavity width for each line
            FP_AMPL: numpy array, the FP line amplitudes
            FP_LL_POS_NEW: numpy array, the corrected wavelengths of the
                           FP lines
            ALL_LINES_2: list of numpy arrays, length = number of orders
                         each numpy array contains gaussian parameters
                         for each found line in that order
\end{pythondocstring}
\end{minipage}

%----------------------------------------------------------------------------------------
\noindent\begin{minipage}{\textwidth}
\subsection{GetE2DSll}

Defined in \spirouTHORCA\path{.spirouTHORCA.get_e2ds_ll}

\begin{pythonbox}
from SpirouDRS import spirouTHORCA
spirouTHORCA.GetE2DSll(p, hdr=None, filename=None, key=None)
spirouTHORCA.spirouTHORCA.get_e2ds_ll(p, hdr=None, filename=None, key=None)
\end{pythonbox}

\begin{pythondocstring}
Get the line list for the e2ds file from "filename" or from calibration
database using hdr (aqctime) and key. Line list is constructed from
fit coefficients stored in keywords:
    'kw_TH_ORD_N', 'kw_TH_LL_D', 'kw_TH_NAXIS1'

:param pp: parameter dictionary, ParamDict containing constants
    Must contain at least:
            log_opt: string, log option, normally the program name
            kw_TH_COEFF_PREFIX: list, the keyword store for the prefix to
                                use to get the TH line list fit coefficients

:param hdr: dictionary or None, the HEADER dictionary with the acquisition
            time in to use in the calibration database to get the filename
            with key=key (or if None key='WAVE_AB')
:param filename: string or None, the file to get the line list from
                 (overrides getting the filename from calibration database)
:param key: string or None, if defined the key in the calibration database
            to get the file from (using the HEADER dictionary to deal with
            calibration database time constraints for duplicated keys.

:return ll: numpy array (1D), the line list values
:return param_ll: numpy array (1d), the line list fit coefficients (used to
                  generate line list - read from file defined)
\end{pythondocstring}
\end{minipage}

%----------------------------------------------------------------------------------------
\noindent\begin{minipage}{\textwidth}
\subsection{GetLampParams}

Defined in \spirouTHORCA\path{.spirouTHORCA.get_lamp_parameters}

\begin{pythonbox}
from SpirouDRS import spirouTHORCA
spirouTHORCA.GetLampParams(p, filename=None, kind=None)
spirouTHORCA.spirouTHORCA.get_lamp_parameters(p, filename=None, kind=None)
\end{pythonbox}

\begin{pythondocstring}
Get lamp parameters from either a specified lamp type="kind" or a filename
or from p['ARG_FILE_NAMES'][0] (if no filename or kind defined)

:param p: parameter dictionary, ParamDict containing constants
    Must contain at least:
        IC_LAMPS: list of strings, the different allowed lamp types
        log_opt: string, log option, normally the program name
:param filename: string or None, the filename to check for the lamp
                 substring in
:param kind: string or None, the lamp type

:return p: parameter dictionary, the updated parameter dictionary
        Adds the following:
            LAMP_TYPE: string, the type of lamp (e.g. UNe or TH)
            IC_LL_LINE_FILE: string, the file name of the line list to use
            IC_CAT_TYPE: string, the line list catalogue type
\end{pythondocstring}
\end{minipage}

%----------------------------------------------------------------------------------------
\noindent\begin{minipage}{\textwidth}
\subsection{JoinOrders}

Defined in \spirouTHORCA\path{.spirouTHORCA.join_orders}

\begin{pythonbox}
from SpirouDRS import spirouTHORCA
spirouTHORCA.JoinOrders
spirouTHORCA.spirouTHORCA.join_orders
\end{pythonbox}

\begin{pythondocstring}
Merge the littrow extrapolated solutions with the fitted line solutions

:param p: parameter dictionary, ParamDict containing constants
    Must contain at least:
        IC_HC_N_ORD_START: int, defines first order solution is calculated
        IC_HC_N_ORD_FINAL: int, defines last order solution is calculated
                            from

:param loc: parameter dictionary, ParamDict containing data
    Must contain at least:
        LL_OUT_2: numpy array (2D), the output wavelengths for each
                  pixel and each order (in the shape of original image)
        DLL_OUT_2: numpy array (2D), the output delta wavelengths for
                   each pixel and each order (in the shape of original
                   image)
        LITTROW_EXTRAP_SOL_2: numpy array (2D),
                          size=([no. orders] by [no. cut points])
                          the wavelength solution at each cut point for
                          each order
        LITTROW_EXTRAP_PARAM_2: numy array (2D),
                          size=([no. orders] by [the fit degree +1])
                          the coefficients of the fits for each cut
                          point for each order
    
:return loc: parameter dictionary, the updated parameter dictionary
    Adds/updates the following:
        LL_FINAL: numpy array, the joined littrow extrapolated and fitted
                  solution wavelengths
        LL_PARAM_FINAL: numpy array, the joined littrow extrapolated and
                        fitted fit coefficients
\end{pythondocstring}
\end{minipage}

%----------------------------------------------------------------------------------------
\noindent\begin{minipage}{\textwidth}
\subsection{SecondGuessSolution}

Defined in \spirouTHORCA\path{.spirouTHORCA.second_guess_at_wave_solution}

\begin{pythonbox}
from SpirouDRS import spirouTHORCA
spirouTHORCA.SecondGuessSolution(p, loc, mode=0)
spirouTHORCA.spirouTHORCA.second_guess_at_wave_solution(p, loc, mode=0)
\end{pythonbox}

\begin{pythondocstring}
Second guess at wave solution, consistency check, using the wavelength
solutions line list

:param p: parameter dictionary, ParamDict containing constants
    Must contain at least:
        IC_LL_SP_MIN: int, minimum wavelength of the catalog
        IC_LL_SP_MAX: int, maximum wavelength of the catalog
        IC_RESOL: int, Resolution of spectrograph
        IC_LL_FREE_SPAN_2: int, window size in sigma unit
        IC_HC_N_ORD_START_2: int, defines first order solution is
                             calculated from
        IC_HC_N_ORD_FINAL_2: int, defines last order solution is
                             calculated from
        IC_HC_T_ORDER_START: int, defines the echelle order of
                            the first e2ds order

:param loc: parameter dictionary, ParamDict containing data
    Must contain at least:
        ECHELLE_ORDERS: numpy array (1D), the echelle order numbers
        HCDATA: numpy array (2D), the image data
        LL_LINE: numpy array, the line list wavelengths from file
        AMPL_LINE: numpy array, the line list amplitudes from file
        LITTROW_EXTRAP_SOL_1: numpy array (2D),
                          size=([no. orders] by [no. cut points])
                          the wavelength solution at each cut point for
                          each order

:param mode: string, if mode="new" uses python to work out gaussian fit
                     if mode="old" uses FORTRAN (testing only) requires
                     compiling of FORTRAN fitgaus into spirouTHORCA dir

:return loc: parameter dictionary, the updated parameter dictionary
        Adds/updates the following:
            ALL_LINES_2: list of numpy arrays, length = number of orders
                         each numpy array contains gaussian parameters
                         for each found line in that order

ALL_LINES_2 definition:
    ALL_LINES_2[row] = [gparams1, gparams2, ..., gparamsN]
                where:
                    gparams[0] = output wavelengths
                    gparams[1] = output sigma (gauss fit width)
                    gparams[2] = output amplitude (gauss fit)
                    gparams[3] = difference in input/output wavelength
                    gparams[4] = input amplitudes
                    gparams[5] = output pixel positions
                    gparams[6] = output pixel sigma width
                                      (gauss fit width in pixels)
                    gparams[7] = output weights for the pixel position
\end{pythondocstring}
\end{minipage}

%----------------------------------------------------------------------------------------


%%%%%%%%%%%%%%%%%%%%%%%%%%%%%%%%%%%%%%%%%%%%%%%%%%%%%%%%
%%
% \section{The spirouUnitTests module}
% \label{ch:the_module:spirouUnitTests}
%%
%%%%%%%%%%%%%%%%%%%%%%%%%%%%%%%%%%%%%%%%%%%%%%%%%%%%%%%%

% % Appendi
% \appendix

% % \addtocontents{toc}{\setlength\cftchapternumwidth{1em}}

% % Appendix A: env_setup.sh
% \input{Chapters/appendix_env_setup_sh}

%%%-------------------------------------------------------------------------------

\backmatter

%%% BIBLIOGRAPHY
%%% -------------------------------------------------------------

% \bibliographystyle{utphysics}
% \bibliography{ref}

\end{document}
